% Options for packages loaded elsewhere
\PassOptionsToPackage{unicode}{hyperref}
\PassOptionsToPackage{hyphens}{url}
\PassOptionsToPackage{dvipsnames,svgnames,x11names}{xcolor}
%
\documentclass[12pt]{article}
\usepackage{amsthm}

% Define theorem-like environments
\newtheorem{definition}{Definition}[section] % numbered by section
\newtheorem{theorem}{Theorem}[section]
\newtheorem{lemma}{Lemma}[section]
\newtheorem{proposition}{Proposition}[section]
\newtheorem{corollary}{Corollary}[section]
\newtheorem{remark}{Remark}[section]

\usepackage{algorithm}
\usepackage{algpseudocode}

\usepackage{amsmath,amssymb}
\usepackage{bm}
\usepackage{tabularx}
\usepackage{iftex}
\ifPDFTeX
  \usepackage[T1]{fontenc}
  \usepackage[utf8]{inputenc}
  \usepackage{textcomp} % provide euro and other symbols
\else % if luatex or xetex
  \usepackage{unicode-math}
  \defaultfontfeatures{Scale=MatchLowercase}
  \defaultfontfeatures[\rmfamily]{Ligatures=TeX,Scale=1}
\fi
\usepackage{lmodern}
\ifPDFTeX\else  
    % xetex/luatex font selection
\fi
% Use upquote if available, for straight quotes in verbatim environments
\IfFileExists{upquote.sty}{\usepackage{upquote}}{}
\IfFileExists{microtype.sty}{% use microtype if available
  \usepackage[]{microtype}
  \UseMicrotypeSet[protrusion]{basicmath} % disable protrusion for tt fonts
}{}
\makeatletter
\@ifundefined{KOMAClassName}{% if non-KOMA class
  \IfFileExists{parskip.sty}{%
    \usepackage{parskip}
  }{% else
    \setlength{\parindent}{0pt}
    \setlength{\parskip}{6pt plus 2pt minus 1pt}}
}{% if KOMA class
  \KOMAoptions{parskip=0}}
\makeatother
\usepackage{xcolor}
\setlength{\emergencystretch}{3em} % prevent overfull lines
\setcounter{secnumdepth}{5}
% Make \paragraph and \subparagraph free-standing
\makeatletter
\ifx\paragraph\undefined\else
  \let\oldparagraph\paragraph
  \renewcommand{\paragraph}{
    \@ifstarcorollary
      \xxxParagraphStar
      \xxxParagraphNoStar
  }
  \newcommand{\xxxParagraphStar}[1]{\oldparagraph*{#1}\mbox{}}
  \newcommand{\xxxParagraphNoStar}[1]{\oldparagraph{#1}\mbox{}}
\fi
\ifx\subparagraph\undefined\else
  \let\oldsubparagraph\subparagraph
  \renewcommand{\subparagraph}{
    \@ifstar
      \xxxSubParagraphStar
      \xxxSubParagraphNoStar
  }
  \newcommand{\xxxSubParagraphStar}[1]{\oldsubparagraph*{#1}\mbox{}}
  \newcommand{\xxxSubParagraphNoStar}[1]{\oldsubparagraph{#1}\mbox{}}
\fi
\makeatother


\providecommand{\tightlist}{%
  \setlength{\itemsep}{0pt}\setlength{\parskip}{0pt}}\usepackage{longtable,booktabs,array}
\usepackage{calc} % for calculating minipage widths
% Correct order of tables after \paragraph or \subparagraph
\usepackage{etoolbox}
\makeatletter
\patchcmd\longtable{\par}{\if@noskipsec\mbox{}\fi\par}{}{}
\makeatother
% Allow footnotes in longtable head/foot
\IfFileExists{footnotehyper.sty}{\usepackage{footnotehyper}}{\usepackage{footnote}}
\makesavenoteenv{longtable}
\usepackage{graphicx}
\makeatletter
\def\maxwidth{\ifdim\Gin@nat@width>\linewidth\linewidth\else\Gin@nat@width\fi}
\def\maxheight{\ifdim\Gin@nat@height>\textheight\textheight\else\Gin@nat@height\fi}
\makeatother
% Scale images if necessary, so that they will not overflow the page
% margins by default, and it is still possible to overwrite the defaults
% using explicit options in \includegraphics[width, height, ...]{}
\setkeys{Gin}{width=\maxwidth,height=\maxheight,keepaspectratio}
% Set default figure placement to htbp
\makeatletter
\def\fps@figure{htbp}
\makeatother

\addtolength{\oddsidemargin}{-.5in}%
\addtolength{\evensidemargin}{-.1in}%
\addtolength{\textwidth}{1in}%
\addtolength{\textheight}{1.7in}%
\addtolength{\topmargin}{-1in}
\makeatletter
\@ifpackageloaded{caption}{}{\usepackage{caption}}
\AtBeginDocument{%
\ifdefined\contentsname
  \renewcommand*\contentsname{Table of contents}
\else
  \newcommand\contentsname{Table of contents}
\fi
\ifdefined\listfigurename
  \renewcommand*\listfigurename{List of Figures}
\else
  \newcommand\listfigurename{List of Figures}
\fi
\ifdefined\listtablename
  \renewcommand*\listtablename{List of Tables}
\else
  \newcommand\listtablename{List of Tables}
\fi
\ifdefined\figurename
  \renewcommand*\figurename{Figure}
\else
  \newcommand\figurename{Figure}
\fi
\ifdefined\tablename
  \renewcommand*\tablename{Table}
\else
  \newcommand\tablename{Table}
\fi
}
\@ifpackageloaded{float}{}{\usepackage{float}}
\floatstyle{ruled}
\@ifundefined{c@chapter}{\newfloat{codelisting}{h}{lop}}{\newfloat{codelisting}{h}{lop}[chapter]}
\floatname{codelisting}{Listing}
\newcommand*\listoflistings{\listof{codelisting}{List of Listings}}
\makeatother
\makeatletter
\makeatother
\makeatletter
\@ifpackageloaded{caption}{}{\usepackage{caption}}
\@ifpackageloaded{subcaption}{}{\usepackage{subcaption}}
\makeatother


\ifLuaTeX
  \usepackage{selnolig}  % disable illegal ligatures
\fi
\usepackage[]{natbib}
\bibliographystyle{agsm}
\usepackage{bookmark}

\IfFileExists{xurl.sty}{\usepackage{xurl}}{} % add URL line breaks if available
\urlstyle{same} % disable monospaced font for URLs
\hypersetup{
  pdftitle={Ranked Set Sampling for High-Breakdown Robust Estimation:
\\ Unified Theory and Finite-Sample Guarantees},
  pdfauthor={Jaafar Almutawa},
  pdfkeywords={High-Breakdown estimators, Ranked Set Sampling, Statistical Efficiency, Contaminated Data Models},
  colorlinks=true,
  linkcolor={blue},
  filecolor={Maroon},
  citecolor={Blue},
  urlcolor={Blue},
  pdfcreator={LaTeX via pandoc}}




\newcommand{\anon}{1}

%set the key \texttt{anon} to ``0'' to hide the authors and acknowledgements,
%  producing the required anonymized version. 
%Set the key \texttt{anon} to ``1'' to produce the manuscript with author details and
% acknowledgments. 

\begin{document}


\def\spacingset#1{\renewcommand{\baselinestretch}%
{#1}\small\normalsize} \spacingset{1}


%%%%%%%%%%%%%%%%%%%%%%%%%%%%%%%%%%%%%%%%%%%%%%%%%%%%%%%%%%%%%%%%%%%%%%%%%%%%%%


\if1\anon
{
  \title{\bf
  Ranked Set Sampling for High-Breakdown Robust Estimation with Unified Theory and Finite-Sample Guarantees
  }
  \author{Jaafar Almutawa \thanks{
I am grateful to Bahrain Polytechnic for providing the facilities and support that enabled this research, and to my colleagues for their encouragement and helpful discussions}\hspace{.2cm}\\
    Department of Mathematics, Bahrain Polytechnic\\
    email: jaafar111@yahoo.com
%    and \\
%    Author 2 \\
%    Department of ZZZ, University of WWW
}
  \maketitle
} \fi

\if0\anon
{
  \bigskip
  \bigskip
  \bigskip
  \begin{center}
    {\LARGE\bf Title}
\end{center}
  \medskip
} \fi

\bigskip
\begin{abstract}
We propose a unified framework integrating Ranked Set Sampling (RSS) with high-breakdown estimators (MCD, MVE, Stahel-Donoho) for robust multivariate estimation. By exploiting RSS's structured subsampling, we preserve maximal breakdown points while accelerating computation. Theoretically, RSS-based estimators retain affine equivariance, and achieve controlled estimation error under $\epsilon$-contamination, with rigorous breakdown proof for RSS-MCD.
Simulations demonstrate substantial efficiency gains (relative efficiency up to 2.8 in univariate settings) without robustness loss, while the RSS ranking overhead remains negligible relative to the cost of the base robust estimator. Multivariate simulations across dimensions $p \in \{4, 10, 50\}$ reveal that RSS gains are concentrated in specific estimator--contamination combinations, with the trimmed mean benefiting most (RE up to 1.32 at $p=4$), and that gains diminish but persist as dimensionality increases.
Crucially, RSS reduces the subsampling complexity from \(O(n^k)\) to \(O(m^c)\) for set size \(m\) and cycles \(c\), enabling scalable robust estimation.
\end{abstract}

\noindent%
{\it Keywords:} High-Breakdown estimators, Ranked Set Sampling, Statistical Efficiency, Contaminated Data Models
\vfill

\newpage
\spacingset{1.0} % DON'T change the spacing!

\section{Introduction}
Robust estimation of multivariate location and scatter is a fundamental challenge in statistical analysis, particularly in the presence of outliers or contamination. Classical estimators, such as the sample mean and covariance matrix, are known to perform optimally under ideal conditions like multivariate normality. However, their extreme sensitivity to even a small fraction of anomalous observations renders them unreliable for real-world data, where deviations from assumptions are common due to measurement errors, data corruption, or inherent variability in complex systems \cite{maronna2019robust}.

To address this challenge, a range of robust estimators have been proposed that trade off some statistical efficiency for significantly enhanced resistance to outliers \cite{maronna2019robust}.
Among these, high-breakdown point estimators—such as the Minimum Covariance Determinant (MCD) \cite{rousseeuw1999fast,hubert2018minimum,
Mayrhofer2025},
Minimum Volume Ellipsoid (MVE) \cite{rousseeuw1999fast},
S-estimators \cite{maronna2019robust}, and the Stahel-Donoho estimator \cite{maronna2019robust}
are widely regarded for their ability to withstand contamination levels approaching $50\%$. These estimators are typically affine equivariant, ensuring that their performance remains invariant under transformations such as scaling, rotation, or translation—an essential property in multivariate settings.

Despite their strong robustness, a recurring limitation of high-breakdown estimators lies in their reduced efficiency under clean data conditions, as well as their substantial computational demands, particularly for large-sample contexts \cite{Hubert2012,Raymaekers2023}. This motivates the need for enhanced strategies that preserve robustness while improving statistical efficiency and practical applicability.

%\textbf{Computational Motivation:}
The combinatorial nature of high-breakdown estimators imposes severe scalability limits \cite{maronna2019robust}.
For MCD with $n=1000$ and $h=500$, exhaustive subsampling requires evaluating $\binom{1000}{500} \approx 10^{299}$ subsets. Even approximate algorithms like FAST-MCD \cite{hubert2018minimum} remain $O(n^2 p)$ in practice. Ranked Set Sampling addresses this by strategically selecting $m$-sized ranked subsets, reducing subsample evaluations to $O(m^c)$ while preserving statistical guarantees.

In this paper, we investigate a principled enhancement of high-breakdown estimators using \emph{Ranked Set Sampling} (RSS)
\cite{McIntyre1952,zhang2024birnbaum}, an intelligent sampling design originally developed for improving estimation efficiency when measurement is expensive but ranking is easy or cheap. In the context of robust multivariate estimation, ranking can be efficiently performed using projection-based depth measures or robust univariate statistics \cite{Hubert2012},
making RSS especially attractive. We propose a unified theoretical and algorithmic framework that integrates RSS into the computation of robust location and scatter estimators. Our approach systematically improves statistical efficiency while preserving both the breakdown point and affine equivariance of classical robust estimators.

The contributions of this paper are threefold:
\begin{enumerate}
    \item We develop a comprehensive methodology for enhancing high-breakdown estimators via RSS, supported by rigorous theoretical analysis of breakdown point preservation, efficiency gains, concentration inequalities, and computational complexity.
    \item We review the foundational high-breakdown estimators—MCD, MVE, S-estimators, and Stahel-Donoho—and summarize their key theoretical properties, robustness guarantees, and computational characteristics.
    \item We validate the proposed framework through theoretical performance bounds and provide algorithmic implementations to guide practical deployment in real-world applications.
\end{enumerate}

%\textbf{Novelty Position:}
Our work presents the first unified integration of RSS with high-breakdown multivariate estimators. Unlike random sampling accelerators that sacrifice breakdown points \cite{maronna2019robust} or resampling methods requiring $O(Bn)$ computations \cite{rousseeuw1999fast}, our method maintains the original estimators' breakdown properties while achieving quantifiable speedups. This bridges a critical gap between robust statistics and efficient computation.

While RSS has been used to improve efficiency in other contexts \cite{yusuf2023novel}, its integration with high-breakdown multivariate estimators remains largely unexplored in the literature. Multivariate ranking via depth measures introduces additional computational cost, which we address through efficient ranking strategies in Section~\ref{sec:methodology}.

The proposed methodology has significant implications for applications in fields where robust multivariate analysis is crucial. These include environmental monitoring (e.g., pollution measurements with irregular sampling) and engineering systems (e.g., sensor fusion under noise and drift). It is also relevant for financial risk modeling (e.g., detection of market anomalies) and biomedical data analysis (e.g., outlier-prone high-throughput measurements). In all these domains, maintaining reliability under data contamination while improving estimator efficiency is of paramount importance.

To the best of our knowledge, this is the first work to formally integrate RSS with high-breakdown multivariate estimators and to provide theoretical guarantees for the combined framework.


The rest of the paper is organized as follows. Section~\ref{sec:theoretical} reviews the theoretical foundations of high-breakdown estimators and summarizes their key robustness and equivariance properties.
Section~\ref{sec:rss} discusses Ranked Set Sampling in the context of robust univariate and multivariate estimation.
Section~\ref{sec:methodology} lays out our proposed methodology, showing how the RSS-enhanced robust estimation framework can be implemented in practice, with attention to both the algorithmic design and practical considerations.
Finally, Section~\ref{sec:simulation} reports the results of simulation studies including univariate experiments, multivariate experiments across dimensions $p \in \{4, 10, 50\}$, and illustrative applications on benchmark datasets evaluating the performance of the proposed approach.

\section{High-Breakdown Estimators: Theory and Computational Considerations}
\label{sec:theoretical}

High-breakdown estimators are indispensable tools in multivariate statistical analysis, particularly when data may be contaminated with outliers. These estimators balance robustness with statistical efficiency and are capable of handling large contamination levels without severe degradation in performance \cite{zhang2023fast}.
This section presents the theoretical foundations and computational properties of several widely used high-breakdown estimators for location and scatter, providing a foundation for the efficiency improvements proposed in later sections through Ranked Set Sampling (RSS), which aims to reduce the computational burden inherent in subset searches or projection-based diagnostics without sacrificing robustness.

Let $T_n$ be an estimator based on a sample $X_n = \{x_1, \ldots, x_n\}$. The finite-sample breakdown point of $T_n$
at $X_n$ is defined as \cite{rousseeuw1999fast}:
\[
\epsilon^*_n(T_n, X_n) = \frac{1}{n} \min \left\{ m : \sup_{X_n'} \| T_n(X_n') - T_n(X_n) \| = \infty \right\},
\]
where $X_n'$ is obtained by replacing $m$ observations in $X_n$ with arbitrary values. In other words, the breakdown point quantifies the smallest proportion of contaminated data that can cause the estimator to yield arbitrarily large errors. The asymptotic breakdown point is the limit of $\epsilon^*_n$ as $n \to \infty$, representing the estimator’s ultimate robustness against increasing sample size.

For affine equivariant estimators, the following theoretical upper bounds apply \cite{maronna2019robust}:
\[
\epsilon^*_n \leq \frac{1}{n} \left\lfloor \frac{n + 1}{2} \right\rfloor \quad \text{(location)}, \quad \epsilon^*_n \leq \frac{1}{n} \left\lfloor \frac{n - p + 1}{2} \right\rfloor \quad \text{(scatter)},
\]
where $p$ is the dimension of the data. These bounds establish the fundamental limits of robustness achievable for affine equivariant estimators. The maximum attainable breakdown point of $50\%$ is considered optimal, as it corresponds to the highest proportion of contamination that can be tolerated without complete loss of estimator reliability~\cite{maronna2019robust}.

High-breakdown estimators (with $\epsilon^* \approx 0.5$)
generally possess the following features \cite{maronna2019robust}:
\begin{itemize}
    \item \textbf{Affine Equivariance:} For any affine transformation $(A \mathbf{x} + \mathbf{b})$, these estimators ensure consistent results across transformations, formally expressed as:
    \[
    \hat{\bm{\mu}}(A\mathbf{X} + \mathbf{b}) = A\hat{\bm{\mu}}(\mathbf{X}) + \mathbf{b}, \quad \hat{\bm{\Sigma}}(A\mathbf{X} + \mathbf{b}) = A\hat{\bm{\Sigma}}(\mathbf{X})A^\top.
    \]
    \item \textbf{Subset or Projection-Based Resistance:} They rely on robust subsets (e.g., MCD, MVE) or projection-based diagnostics (e.g., Stahel-Donoho) to identify and downweight or exclude outliers.
    \item \textbf{Distribution-Free Robustness:} These estimators do not require specific distributional assumptions on the outlier-generating mechanism for their breakdown point guarantees. However, efficiency results (such as asymptotic relative efficiency) may depend on the underlying distribution and the ranking error model; see Section~3 for a detailed discussion.
    \item \textbf{Efficiency-Robustness Trade-off:} While achieving high breakdown points, these estimators typically exhibit reduced statistical efficiency under clean data conditions compared to classical maximum likelihood estimators. Efficiency is often measured as asymptotic relative efficiency (ARE) under the multivariate normal model.
\end{itemize}

Robust multivariate estimation is essential in practice since classical estimators are sensitive to contamination. High-breakdown estimators mitigate this issue by tolerating a significant portion of arbitrary contamination without catastrophic performance degradation. However, their reduced efficiency under ideal conditions and substantial computational demands present opportunities for improvement through intelligent sampling designs such as Ranked Set Sampling. To provide a detailed understanding of the available tools, we now present four foundational high-breakdown estimators, evaluating each in terms of robustness, efficiency, computational practicality, and specific scenarios where they excel.

\subsection{Minimum Covariance Determinant (MCD)}
The MCD estimator~\cite{rousseeuw1999fast} identifies a subset $H \subset \{1, \ldots, n\}$ of size $h \approx n/2$ that minimizes the determinant of the sample covariance matrix:
\[
\det(\hat{\bm{\Sigma}}_H) = \det\left( \frac{1}{h - 1} \sum_{i \in H} (\mathbf{x}_i - \bar{\mathbf{x}}_H)(\mathbf{x}_i - \bar{\mathbf{x}}_H)^\top \right),
\]
where $\bar{\mathbf{x}}_H = \frac{1}{h} \sum_{i \in H} \mathbf{x}_i$. The mean and covariance of this optimal subset serve as the raw MCD estimates. The MCD estimator is widely used for robust estimation of multivariate location and scatter.

A reweighting step based on Mahalanobis distances from the raw MCD estimates is typically used to improve efficiency. The FAST-MCD algorithm~\cite{rousseeuw1999fast} accelerates computation, reducing complexity from $O(n^3)$ to $O(n^2p)$ via intelligent subsampling and iterative refinement.

The breakdown point for $h = \left\lfloor \frac{n + p + 1}{2} \right\rfloor$ is
\[
\epsilon^*_n = \frac{n - h + 1}{n} \approx \frac{n - p + 1}{2n} \rightarrow 0.5 \quad \text{as } n \to \infty.
\]
Thus, MCD can tolerate nearly $50\%$ contamination. Its asymptotic efficiency at the normal model is about $85\%$ after reweighting~\cite{maronna2019robust}. However, MCD requires $n > 2p$ for the estimator to be well-defined, and its efficiency is lower than that of the sample mean/covariance under clean data.

%\textbf{Practical limitations and the value of RSS:}
In practice, the MCD can quickly become too demanding to compute when the number of variables is large compared to the sample size, even with faster approximations such as FAST-MCD. Its performance also depends on how the subset size $h$ is chosen and on the stability of the reweighting step, which can make results somewhat sensitive. Ranked Set Sampling (RSS) provides a way to narrow down the search more intelligently by ranking observations before selecting subsets, which helps cut down the computational cost while still keeping the estimator robust.
\subsection{S-Estimators}
S-estimators~\cite{maronna2019robust} are defined by minimizing a robust scale $\sigma$ subject to a constraint:
\[
\min_{\bm{\mu}, \bm{\Sigma}} \sigma \quad \text{subject to} \quad \frac{1}{n} \sum_{i=1}^n \rho\left( \frac{d_i^2}{\sigma^2} \right) = \delta,
\]
where $d_i$ is the Mahalanobis distance, $\rho$ is a bounded loss function (e.g., Tukey's biweight), and $\delta$ is a constant for consistency.
S-estimators achieve breakdown points approaching $0.5$ with suitable $\rho$ and $\delta$. Their efficiency is tunable: with appropriate $\rho$, ARE can exceed $90\%$ under normality~\cite{maronna2019robust}. However, computation is intensive: $O(n^2p^2)$ per iteration.

%\textbf{Practical limitations and RSS integration:}
High-breakdown estimators often come with practical hurdles: they can be sensitive to tuning choices, and convergence may slow considerably as the sample size or dimensionality increases. By incorporating Ranked Set Sampling (RSS), the search for candidate subsets or projections can be made more focused, reducing the number of evaluations needed in the optimization and easing the overall computational burden.

\subsection{Minimum Volume Ellipsoid (MVE)}
The MVE estimator~\cite{rousseeuw1985multivariate} finds the smallest ellipsoid (minimum determinant of $\bm{\Sigma}$) containing at least $h$ data points:
\[
\#\left\{ i : (\mathbf{x}_i - \bm{\mu})^\top \bm{\Sigma}^{-1} (\mathbf{x}_i - \bm{\mu}) \leq c \right\} \geq h
\]
with $c$ chosen for consistency.
Asymptotically, the breakdown point approaches $0.5$. However, MVE has low efficiency (typically $<50\%$ at normality) and is computationally more demanding than MCD ($O(n^3)$ for naive search).

%\textbf{Practical limitations and RSS integration:}
Despite their strong robustness, high-breakdown estimators often suffer from heavy computational costs, which limits their use in large-scale applications. Ranked Set Sampling (RSS) offers a way forward by guiding the selection of candidate ellipsoid subsets in a more efficient manner, thereby reducing the search space and making these estimators more practical for high-dimensional or large-sample problems.
\subsection{Stahel-Donoho Estimator}
The Stahel-Donoho estimator~\cite{maronna2019robust} quantifies outlyingness as
\[
o(\mathbf{x}_i) = \sup_{\|\mathbf{u}\| = 1} \frac{|\mathbf{u}^\top \mathbf{x}_i - \text{med}(\mathbf{u}^\top X)|}{\text{MAD}(\mathbf{u}^\top X)},
\]
using the direction $\mathbf{u}$ where the standardized deviation is maximized. In practice, random projections approximate the supremum, yielding $O(n^2p)$ complexity.
Breakdown point approaches $0.5$ with enough projections. Efficiency is moderate and can be improved with reweighting.

%\textbf{Practical limitations and RSS integration:}
Projection-based estimators often depend heavily on both the number and the choice of projection directions, and their computational cost can quickly escalate as $n$ and $p$ grow. By integrating Ranked Set Sampling (RSS), projection directions or initial subsets for scale estimation can be selected in a more targeted way, improving efficiency without sacrificing robustness.

\begin{table}[H]
\centering
\caption{Comparison of High-Breakdown Estimators}
\label{tab:estimator_comparison}
\begin{tabularx}{\textwidth}{lXXXX}
\toprule
\textbf{Estimator} & \textbf{Breakdown Point} & \textbf{Efficiency (Normal)} & \textbf{Computational Cost} & \textbf{RSS Integration Potential} \\
\midrule
MCD           & $\approx 0.5$ & Moderate ($\sim$85\%)  & $O(n^2p)$ (FAST-MCD)      & Subset selection            \\
S-Estimator   & $\approx 0.5$ & High ($>90\%$) (tunable) & $O(n^2p^2)$ per iteration & Subset/projection selection \\
MVE           & $\approx 0.5$ & Low ($<50\%$)           & $O(n^3)$                  & Subset selection            \\
Stahel-Donoho & $\approx 0.5$ & Moderate                & $O(n^2p)$ (random proj.)  & Projection/subset selection \\
\bottomrule
\end{tabularx}
\end{table}

In summary, high-breakdown estimators play a critical role in ensuring robustness against data contamination, each offering unique strengths and trade-offs. Their combination with efficient data collection methods---as detailed in subsequent sections---forms a promising strategy for enhancing both robustness and statistical performance in practical applications.

\section{Ranked Set Sampling for Robust Multivariate Estimation}\label{sec:rss}

Ranked Set Sampling (RSS), first proposed by McIntyre \cite{McIntyre1952},
is a probability sampling technique that utilizes auxiliary ranking information to improve estimation efficiency while maintaining unbiasedness. By systematically prioritizing high-information-density regions within a population, RSS achieves variance reduction without increasing sample size \cite{Alvandi2023}.
Integrating RSS with robust estimation thus offers dual benefits: enhanced efficiency through optimized sampling and sustained resistance to outliers, effectively addressing the classical efficiency-breakdown trade-off inherent in robust statistics.
RSS is particularly advantageous when direct measurement is expensive or invasive, but ranking can be done cheaply or non-destructively. This makes RSS highly relevant for robust estimation, where both efficiency and resistance to outliers are critical.
The integration of RSS with robust estimation allows for more efficient estimators that retain the robustness properties of traditional robust methods, effectively addressing the trade-off between efficiency and breakdown.
\subsection{Ranked Set Sampling: Fundamentals and Motivation}

Ranked Set Sampling is a cost-effective sampling design that improves estimation efficiency by leveraging ranking information obtained from inexpensive or non-destructive assessments. Given a set size $m$ and $k$ independent cycles, the RSS protocol comprises three systematic steps \cite{Wolfe2010}:

\begin{enumerate}
    \item \textbf{Selection}: Independently draw $k$ simple random samples, each consisting of $m$ units.
    \item \textbf{Ranking}: Within each sample, rank the $m$ units based on an auxiliary variable or expert judgment, typically without precise measurement.
    \item \textbf{Measurement}: From the $i$-th sample, measure only the unit ranked $r_i$.
\end{enumerate}
The ranking step is typically performed using a concomitant variable or visual assessment, rather than the actual measurement of interest. This allows for efficient sampling even when precise measurement is costly or time-consuming.
This process yields an RSS sample of the form $\{X_{(r_1)1}, X_{(r_2)2}, \ldots, X_{(r_k)k}\}$, where $X_{(r_i)i}$ denotes the observation with rank $r_i$ in the $i$-th set. This structure ensures each observation contributes information from a diverse spread across the range of the underlying variable.

A widely used special case is the \textit{Balanced RSS} (BRSS), where $k = m$ and $r_i = i$ for all $i = 1, \ldots, m$, ensuring that each rank position contributes exactly once to the final sample: $\{X_{(1)1}, X_{(2)2}, \ldots, X_{(m)m}\}$. Under standard distributional assumptions, the $i$-th order statistic in a sample of size $m$ follows the known density function \cite{ChenStasnyWolfe2006}:
\[
f_{(i)}(x) = \frac{m!}{(i-1)!(m-i)!} [F(x)]^{i-1} [1-F(x)]^{m-i} f(x),
\]
where $f(x)$ and $F(x)$ denote the probability density and cumulative distribution functions, respectively.
In practice, unbalanced RSS—where some ranks are sampled more frequently than others—can be used to further improve precision, especially for skewed distributions.

The RSS sample mean
\[
\bar{X}_{RSS} = \frac{1}{k} \sum_{i=1}^{k} X_{(r_i)i},
\]
remains an unbiased estimator of the population mean $\mu$, irrespective of the accuracy of the ranking mechanism. For perfectly ranked BRSS, its variance is given by \cite{ChenStasnyWolfe2006}:
\[
\text{Var}(\bar{X}_{RSS}) = \frac{1}{m^2} \sum_{i=1}^{m} \text{Var}(X_{(i)}) - \frac{1}{m^2} \sum_{i=1}^{m} \left[ E(X_{(i)}) - \mu \right]^2,
\]
where $\text{Var}(X_{(i)})$ depends on the underlying distribution. Importantly, this variance is guaranteed to be no greater than the variance of the sample mean under simple random sampling (SRS), highlighting one of the core advantages of RSS in improving estimation precision. The relative efficiency of RSS compared to SRS is defined as
\begin{equation}
\label{eq:re_general}
\text{RE}_{\text{RSS}} := \frac{\text{Var}(\bar{X}_{\text{SRS}})}{\text{Var}(\bar{X}_{\text{RSS}})} \geq 1,
\end{equation}
where $\bar{X}_{\text{SRS}}$ is the simple random sample mean and $\bar{X}_{\text{RSS}}$ the ranked set sample mean.

In practice, perfect ranking is rarely attainable, making the impact of ranking errors a central concern in the efficiency of RSS-based estimators. The concept of ranking quality provides a quantitative basis to assess this impact and predict expected gains under real-world conditions.
Even with imperfect ranking, RSS generally leads to more efficient estimation than simple random sampling, as long as the ranking is better than random. If ranking is no better than random, RSS reduces to SRS.
\begin{definition}[Ranking Quality]
\label{def:ranking_quality}
The quality of ranking is quantified by the Spearman rank correlation \cite{taconeli2023,taconeli2023dual}:
\begin{equation}
\rho = \text{Corr}(R, \tilde{R}),
\end{equation}
where $R$ denotes the true ranks and $\tilde{R}$ the observed ranks. This yields:
\begin{enumerate}
    \item \textbf{Perfect Ranking:} $\rho = 1$
    \item \textbf{Random Ranking:} $\rho = 0$
    \item \textbf{Imperfect Ranking:} $0 < \rho < 1$
\end{enumerate}
\end{definition}

The relative efficiency $\text{RE}(\rho)$ quantifies the gain in precision as a function of ranking quality $\rho$, normalized so that $\text{RE}(0) = 1$ (random ranking) and increasing smoothly with $\rho$, reaching its maximum at $\rho=1$ (perfect ranking). This relationship provides a direct link between ranking performance and sampling efficiency, serving as a practical guideline for evaluating RSS under real-world imperfections.

The ranking quality parameter $\rho$ fundamentally determines the efficiency gain achievable through Ranked Set Sampling. The relationship between $\rho$ and statistical efficiency is well-characterized by a quadratic dependence, reflecting how information loss in the ranking process proportionally reduces the advantages of RSS. This quadratic behavior holds under the following key conditions: (1) ranking errors are independent across sets, (2) the ranking mechanism is consistent across cycles, and (3) the population distribution remains stationary during sampling.

\begin{theorem}[RSS Efficiency Under Imperfect Ranking]
\label{thm:rss_efficiency}
Let $\Delta = \mathrm{RE}(1) - 1 > 0$, where $\mathrm{RE}(1)$ is the relative efficiency under perfect ranking. Suppose ranking errors are independent across sets with constant marginal variances. Then the relative efficiency of RSS compared to SRS is given by
\begin{equation}\label{eq:re_rho_exact}
\mathrm{RE}(\rho) = \frac{1 + \Delta}{1 + (1 - \rho)\Delta},
\end{equation}
where $\rho$ is the ranking correlation (Definition~\ref{def:ranking_quality}). Moreover, $\mathrm{RE}(\rho)$ satisfies the lower bound
\begin{equation}\label{eq:re_rho_bound}
\mathrm{RE}(\rho) \geq 1 + \rho^2 \cdot \Delta,
\end{equation}
with equality at $\rho = 0$ and $\rho = 1$. In particular:
\begin{enumerate}
\item $\mathrm{RE}(\rho) \geq 1$ for all $\rho \in [0,1]$.
\item $\mathrm{RE}(\rho)$ is strictly increasing in $\rho$.
\end{enumerate}
\end{theorem}

\begin{proof}
Decompose the RSS mean as
$\bar{X}_{\mathrm{RSS}} = \bar{X}_{\mathrm{perfect}} + \bar{X}_{\mathrm{error}}$,
where $\bar{X}_{\mathrm{perfect}}$ is the idealized RSS mean under perfect ranking and $\bar{X}_{\mathrm{error}}$ captures the effect of ranking imperfections. Assuming ranking errors are independent across sets with constant marginal variances, variance decomposition gives
\[
\mathrm{Var}(\bar{X}_{\mathrm{RSS}}) = \mathrm{Var}(\bar{X}_{\mathrm{perfect}}) + (1 - \rho)\bigl[\mathrm{Var}(\bar{X}_{\mathrm{SRS}}) - \mathrm{Var}(\bar{X}_{\mathrm{perfect}})\bigr].
\]
Writing $\sigma^2_{\mathrm{SRS}} = \mathrm{Var}(\bar{X}_{\mathrm{SRS}})$, $\sigma^2_{\mathrm{P}} = \mathrm{Var}(\bar{X}_{\mathrm{perfect}})$, and $\Delta = \mathrm{RE}(1) - 1 = (\sigma^2_{\mathrm{SRS}} - \sigma^2_{\mathrm{P}})/\sigma^2_{\mathrm{P}}$, we obtain
\[
\mathrm{Var}(\bar{X}_{\mathrm{RSS}}) = \sigma^2_{\mathrm{P}}\bigl[1 + (1 - \rho)\Delta\bigr],
\]
so that
\[
\mathrm{RE}(\rho) = \frac{\sigma^2_{\mathrm{SRS}}}{\mathrm{Var}(\bar{X}_{\mathrm{RSS}})} = \frac{\sigma^2_{\mathrm{P}}(1 + \Delta)}{\sigma^2_{\mathrm{P}}[1 + (1 - \rho)\Delta]} = \frac{1 + \Delta}{1 + (1 - \rho)\Delta},
\]
establishing~\eqref{eq:re_rho_exact}. For the lower bound~\eqref{eq:re_rho_bound}, define $g(\rho) = (1+\Delta)/[1+(1-\rho)\Delta]$ and $h(\rho) = 1 + \rho^2\Delta$. We verify: $g(0) = h(0) = 1$, $g(1) = h(1) = 1 + \Delta$, and
\[
g''(\rho) = \frac{2\Delta^2(1+\Delta)}{[1+(1-\rho)\Delta]^3} > 0,
\]
so $g$ is convex on $[0,1]$. Since $h$ is the secant line connecting $(0, g(0))$ and $(1, g(1))$ adjusted by the quadratic $\rho^2\Delta$, we verify the inequality directly: $g(\rho) - h(\rho) = g(\rho) - 1 - \rho^2\Delta \geq 0$ follows from the convexity of $g$ and the fact that both sides agree at the endpoints, with $g$ curving above $h$ in the interior. Formally, for $\rho \in (0,1)$:
\[
\frac{1+\Delta}{1+(1-\rho)\Delta} \geq 1 + \rho^2\Delta,
\]
which can be verified by cross-multiplying and simplifying to $\rho(1-\rho)\Delta^2 \geq 0$, which holds for all $\rho \in [0,1]$ and $\Delta \geq 0$.
\end{proof}

The closed-form expression for $\text{RE}(\rho)$ highlights its explicit dependence on the ranking correlation $\rho$, yielding several practically relevant properties, including monotonicity and guaranteed improvement over random ranking.

\begin{corollary}
The function $\text{RE}(\rho)$ is strictly increasing in $\rho$ over $[0,1]$ and satisfies the midpoint inequality:
\[
\text{RE}(0.5) \geq 1 + \frac{\text{RE}(1) - 1}{4}.
\]
\end{corollary}

While the derived lower bound holds generally, the exact quadratic form $\text{RE}(\rho) = 1 + \rho^2 (\text{RE}(1) - 1)$ is valid only under specific assumptions. The next remark delineates when this equality is achieved exactly and when it should be interpreted as a conservative approximation.

\begin{remark}
The exact efficiency is given by the closed-form expression~\eqref{eq:re_rho_exact}, which holds under the stated assumptions of independent ranking errors with constant marginal variances. The quadratic lower bound~\eqref{eq:re_rho_bound} is distribution-free in the sense that it holds regardless of the form of the ranking error distribution; it requires only that the ranking correlation $\rho$ is well-defined. The bound is tight at the endpoints $\rho \in \{0, 1\}$ and conservative in the interior, with the gap bounded by $\rho(1-\rho)\Delta^2/[1+(1-\rho)\Delta]$. We emphasize that the breakdown point guarantees of the underlying robust estimators (Section~2) remain distribution-free; it is only the \emph{efficiency} quantification that involves the ranking error structure.
\end{remark}

This section has established the fundamental properties and efficiency gains of RSS, especially in the context of robust estimation. The next section will extend these ideas to multivariate settings and discuss robust ranking functions in more detail.

\subsection{Multivariate Extensions of RSS}
Extending Ranked Set Sampling (RSS) to multivariate settings enables enhanced efficiency in estimating multivariate location and scatter, particularly within the framework of robust statistics. However, multivariate RSS (MRSS) requires careful design of ranking mechanisms that remain informative and resistant to contamination in high-dimensional spaces \cite{taconeli2023}.
\begin{definition}[Multivariate Ranked Set Sampling (MRSS)]
Let $\mathbf{X} = (X_1, X_2, \ldots, X_p)^T$ be a $p$-dimensional random vector from a multivariate population. The MRSS procedure proceeds as follows:
\begin{enumerate}
    \item Independently draw $k$ simple random samples, each of size $m$.
    \item Within each sample, rank the $m$ observations using a multivariate ranking function $R(\mathbf{x})$.
    \item Select and measure one observation per sample based on its rank position (e.g., balanced RSS or general rank scheme).
\end{enumerate}
\end{definition}

The choice of ranking function $R(\mathbf{x})$ is critical for MRSS performance. Common strategies include:
\begin{itemize}
    \item \textbf{Marginal Ranking:} Ranking based on a single variable $X_j$.
    \item \textbf{Aggregate Ranking:} Using a linear combination $R(\mathbf{x}) = \sum_{j=1}^{p} w_j X_j$, where weights $w_j$ incorporate prior knowledge.
    \item \textbf{Distance-Based Ranking:} Ranking by the Euclidean distance $\|\mathbf{x} - \mathbf{c}\|$ to a fixed center $\mathbf{c}$.
    \item \textbf{Depth-Based Ranking:} Employing statistical depth functions (e.g., Mahalanobis or Tukey's halfspace depth).
\end{itemize}
\begin{table}[H]
\centering
\caption{Comparison of Multivariate Ranking Functions}
\label{tab:ranking_comparison}
\begin{tabular}{lccc}
\toprule
\textbf{Ranking Function} & \textbf{Robustness} & \textbf{Affine-Invariant} & \textbf{Cost} \\
\midrule
Marginal Ranking & Low & No & $O(n)$ \\
Aggregate Ranking & Medium & No & $O(n)$ \\
Distance-Based & Medium & No & $O(np)$ \\
Depth-Based & High & Yes & $O(np^2)$ \\
\bottomrule
\end{tabular}
\end{table}
For example, Tukey's halfspace depth: $D(\mathbf{x}) = \inf_{\mathbf{u}} P(\mathbf{u}^T\mathbf{X} \leq \mathbf{u}^T\mathbf{x})$, and Mahalanobis depth: $D(\mathbf{x}) = [1 + (\mathbf{x} - \bm{\mu})^T \bm{\Sigma}^{-1} (\mathbf{x} - \bm{\mu})]^{-1}$.

Marginal and aggregate rankings are computationally simple but may fail to capture joint distributional structure in high dimensions. Depth-based criteria are generally more robust to outliers and better suited for robust multivariate estimation, especially when the ranking function leverages high-breakdown estimators (e.g., robust covariance matrices) \cite{taconeli2023}.

Depth-based ranking typically costs $O(p^3 + np^2)$ per set due to covariance estimation, versus $O(np)$ for distance-based methods.
\begin{theorem}[Componentwise Relative Efficiency in MRSS]
\label{thm:mrss_efficiency}
Let $\mathbf{X} = (X_1, \ldots, X_p)^T$ follow a continuous multivariate distribution, and let $R(\mathbf{X})$ be a scalar ranking function applied in MRSS with set size $m$. For estimating the mean $\mu_j = \mathbb{E}[X_j]$ of component $X_j$, the relative efficiency of MRSS compared to simple random sampling (SRS) is approximately:
\[
\text{RE}_j \approx 1 + \frac{\text{Cov}^2(X_j, R(\mathbf{X}))}{\text{Var}(X_j) \cdot \text{Var}(R(\mathbf{X}))} \cdot C(m),
\]
where $C(m)$ is a design-specific function of set size, satisfying $C(m) \approx \frac{m+1}{2}$ under symmetric distributions and perfect ranking. This quantifies how the informativeness of the ranking, captured by squared correlation, contributes to efficiency gains in estimating $\mu_j$.
\end{theorem}
This approximation assumes $R(\mathbf{X})$ is a consistent estimator of centrality and $C(m)$ is derived from order statistic variances.
\begin{proof}
Let $X_j^{(r)}$ denote the $j$-th component of the observation selected at rank $r$ according to $R(\mathbf{X})$, and define the MRSS estimator over $k$ cycles as:
\[
\bar{X}_{\text{MRSS},j} = \frac{1}{k} \sum_{i=1}^k X_j^{(r_i)}.
\]
Under SRS, $\text{Var}(\bar{X}_{\text{SRS},j}) = \text{Var}(X_j)/k$. With MRSS, when $R(\mathbf{X})$ positively correlates with $X_j$, selecting observations using ranked order statistics reduces variance. Approximating the MRSS variance yields:
\[
\text{Var}(\bar{X}_{\text{MRSS},j}) \approx \frac{1}{k} \left[ \text{Var}(X_j) \left(1 - \rho_j^2 \left(1 - \frac{1}{C(m)} \right) \right) \right],
\]
where $\rho_j = \text{Corr}(X_j, R(\mathbf{X}))$.

Then:
\[
\text{RE}_j = \frac{\text{Var}(\bar{X}_{\text{SRS},j})}{\text{Var}(\bar{X}_{\text{MRSS},j})} \approx \frac{1}{1 - \rho_j^2 \left(1 - \frac{1}{C(m)} \right)}.
\]
For small $\rho_j^2$, expanding gives:
\[
\text{RE}_j \approx 1 + \rho_j^2 \cdot \left(\frac{C(m) - 1}{C(m)}\right),
\]
and noting that
\[
\rho_j^2 = \frac{\text{Cov}^2(X_j, R(\mathbf{X}))}{\text{Var}(X_j) \cdot \text{Var}(R(\mathbf{X}))},
\]
we obtain:
\[
\text{RE}_j \approx 1 + \frac{\text{Cov}^2(X_j, R(\mathbf{X}))}{\text{Var}(X_j) \cdot \text{Var}(R(\mathbf{X}))} \cdot \left(\frac{C(m) - 1}{C(m)}\right).
\]
When $C(m)$ is moderate/large (e.g., $m \geq 5$), $\frac{C(m) - 1}{C(m)} \approx 1$, giving:
\[
\text{RE}_j \approx 1 + \frac{\text{Cov}^2(X_j, R(\mathbf{X}))}{\text{Var}(X_j) \cdot \text{Var}(R(\mathbf{X}))} \cdot C(m),
\]
which completes the proof.
\end{proof}
For $m=5$, $C(m) \approx 3$, meaning even moderate ranking correlation ($\rho_j=0.7$) yields $~50\%$ efficiency gain.
\subsection{RSS for Robust Estimation: Theoretical Framework}

Integrating RSS with robust estimation requires understanding how the sampling design interacts with estimator robustness. While RSS enhances efficiency through informed sampling, robust estimators aim to retain validity under contamination. Their combination demands that ranking mechanisms not compromise robustness \cite{ChenStasnyWolfe2006}.

\begin{definition}[Robust Ranking Function]
A ranking function $R(\mathbf{x})$ in RSS is \emph{robust} if it satisfies:
\begin{enumerate}
    \item \textbf{Breakdown Point Preservation:} Ranking does not reduce the estimator’s breakdown point.
    \item \textbf{Affine Equivariance/Invariance:} Ranking either transforms predictably under affine changes:
    \[
    R(A\mathbf{x} + \mathbf{b}) = g(A, \mathbf{b}) \cdot R(\mathbf{x}),
    \]
    for some function $g$ and full-rank matrix $A$, vector $\mathbf{b}$.
    \item \textbf{Outlier Resistance:} Ranking is insensitive to extremes by relying on robust summaries (e.g., depth functions, robust distances).
\end{enumerate}
\end{definition}

The rationale for robust ranking is to maintain estimator properties. Traditional robust estimators assume simple random sampling; RSS selects observations informed by ranking, which, if sensitive to contamination, could amplify outlier effects. Robust rankings thus preserve representativeness even under contamination \cite{Hubert2012}.

For high-breakdown estimators like MCD, affine equivariance of $R$ ensures the combined estimator retains equivariance.
\begin{theorem}[RSS–Robust Estimator Compatibility]
\label{thm:rss_robust_compatibility}
Let $T_n$ be a robust estimator with breakdown point $\epsilon^*$ under SRS of size $n$, and $T_{RSS}$ the same estimator applied to RSS samples with a robust ranking $R(\mathbf{x})$. Assuming contamination is independent of ranking:
\begin{enumerate}
    \item If $T_n$ is unbiased, so is $T_{RSS}$.
    \item If $T_n$ is consistent as $n \to \infty$, so is $T_{RSS}$.
    \item $\epsilon^*_{RSS} \geq \epsilon^*$.
    \item If $R(\mathbf{x})$ positively correlates with the target parameter, $\text{Var}(T_{RSS}) \leq \text{Var}(T_n)$.
\end{enumerate}
\end{theorem}
Part (iii) is crucial—it means RSS doesn't degrade the robustness of high-breakdown estimators like MCD or S-estimators.
\begin{proof}
(i) Since $\mathbb{E}[X^{(r)}]=\mathbb{E}[X]$ when $R$ is robust and independent of contamination, linear or symmetric estimators preserve unbiasedness.

(ii) MRSS distributions converge weakly to the population under regularity, so consistency follows.

(iii) If contamination is independent of ranking, robust $R$ avoids systematic selection of contaminated points, preserving $\epsilon^*$.

(iv) Positive correlation between $R$ and the parameter reduces conditional variance via variance decomposition, giving $\text{Var}(T_{RSS}) \leq \text{Var}(T_n)$.
\end{proof}
The independence assumption may fail if outliers affect ranking; validate via sensitivity analysis in Section 5.
\subsection{Asymptotic Properties of RSS}

The asymptotic behavior of RSS estimators is key for inference
\cite{ChenStasnyWolfe2006, maronna2019robust}, or \cite{stokes1977efficiency}.

\begin{theorem}[Asymptotic Normality of RSS Estimators]
\label{thm:rss_asymptotic}
Let $\bar{\mathbf{X}}_{RSS} = \frac{1}{n} \sum_{i=1}^{n} \mathbf{X}_{(r_i)i}$ be the RSS sample mean vector. Under standard regularity (finite second moments, Lindeberg condition):
\[
\sqrt{n}(\bar{\mathbf{X}}_{RSS} - \boldsymbol{\mu}) \xrightarrow{d} \mathcal{N}(\mathbf{0}, \boldsymbol{\Sigma}_{RSS}),
\]
where $\boldsymbol{\Sigma}_{RSS}$ is the asymptotic covariance, satisfying $\boldsymbol{\Sigma}_{RSS} \preceq \boldsymbol{\Sigma}_{SRS}$ under perfect ranking.
\end{theorem}

\begin{proof}
Each $\mathbf{X}_{(r_i)i}$ has finite moments and independent across RSS cycles. The Lindeberg-Feller CLT implies asymptotic normality. In balanced designs, $\mathbb{E}[\bar{\mathbf{X}}_{RSS}]=\boldsymbol{\mu}$ and asymptotic variance is:
\[
\boldsymbol{\Sigma}_{RSS} = \frac{1}{m} \sum_{r=1}^{m} \boldsymbol{\Sigma}_r,
\]
with $\boldsymbol{\Sigma}_r$ the covariance of the $r$-th order statistic. Order statistics reduce variance relative to unordered sampling in symmetric distributions, so $\boldsymbol{\Sigma}_{RSS} \preceq \boldsymbol{\Sigma}_{SRS}$.
\end{proof}
This result justifies the use of standard normal-based confidence intervals and hypothesis tests for RSS estimators.
\begin{corollary}[Confidence Intervals for RSS Estimators]
\label{cor:rss_ci}
Under Theorem~\ref{thm:rss_asymptotic}, an approximate $(1-\alpha)$ confidence interval for component $\mu_j$ of $\boldsymbol{\mu}$ is:
\[
\bar{X}_{j,RSS} \pm z_{\alpha/2} \sqrt{\frac{\hat{\sigma}^2_{j,RSS}}{n}},
\]
where $\hat{\sigma}^2_{j,RSS}$ is a consistent estimator of $\text{Var}(\bar{X}_{j,RSS})$.
\end{corollary}
In practice, $\hat{\sigma}^2_{j,RSS}$ can be estimated using the sample variance of the observed order statistics across cycles.
\subsection{RSS for Robust Multivariate Estimation: Advanced Framework}
\label{subsec:advanced_framework}

In extending Ranked Set Sampling (RSS) to robust multivariate estimation, it is essential to ensure that the sampling design complements the robustness characteristics of the estimator. Specifically, high-breakdown estimators such as the Minimum Covariance Determinant (MCD) or S-estimators can benefit from variance reduction through RSS, but only if the ranking procedure
itself is robust to contamination \cite{taconeli2023}.

\begin{definition}[Robust-Compatible RSS Design]
\label{def:robust_rss}
An RSS design is \emph{robust-compatible} if:
\begin{enumerate}
    \item The ranking function $R(\mathbf{x})$ has a bounded influence function.
    \item The breakdown point of the ranking mechanism, $\epsilon^*_R$, satisfies $\epsilon^*_R \geq \epsilon^*_T$, where $\epsilon^*_T$ is the breakdown point of the estimator.
    \item The selected observations maintain sufficient data diversity for robust estimation (e.g., through spatial or distributional coverage).
\end{enumerate}
\end{definition}
For example, using the Mahalanobis distance computed from the MCD estimator as the ranking function yields a robust-compatible RSS design.

The definition of robust-compatible RSS design ensures that the ranking mechanism preserves the robustness properties of the estimator. Building on the asymptotic normality of the RSS sample mean established earlier, we now extend results to general robust multivariate estimators, aiming to characterize their limiting distribution and variance structure. This shows that efficiency gains can be achieved without sacrificing robustness.
{\sloppy
\begin{theorem}[Asymptotic Distribution of Robust Estimators Under RSS]
\label{thm:asymp_robust_rss}
Let $T$ be a multivariate robust estimator with a bounded influence function satisfying $\int \mathrm{IF}^2(\mathbf{x}; T, F)\, dF(\mathbf{x}) < \infty$ (e.g., M-estimators with bounded $\psi$-functions, S-estimators with bounded $\rho$-functions, or the reweighted MCD). Let $T_{\mathrm{RSS}}$ be the same estimator applied to an RSS sample of size $n$, constructed using a robust-compatible ranking function. Then, under regularity conditions ensuring asymptotic linearity of $T$ (see, e.g., Maronna et al.\ 2019, Chapter~6),
\begin{equation}
\sqrt{n}(T_{RSS} - \theta) \xrightarrow{d} \mathcal{N}(0, V_{RSS}(T, F)),
\end{equation}
where the asymptotic variance $V_{RSS}(T, F)$ is given by
\begin{equation}
V_{RSS}(T, F) = \sum_{r=1}^{m} \pi_r \int \mathrm{IF}^2(\mathbf{x}; T, F) f_{(r)}(\mathbf{x}) \, d\mathbf{x},
\end{equation}
with $f_{(r)}$ denoting the density of the $r$-th ranked observation in each RSS cycle and $\pi_r = 1/m$ for balanced designs. Moreover,
\begin{equation}
V_{RSS}(T, F) \leq V_{SRS}(T, F),
\end{equation}
in the sense of positive semidefiniteness, under perfect or near-perfect ranking.
\end{theorem}
}
\begin{proof}
Since $T$ admits an influence function, its asymptotic variance under i.i.d. sampling is
\begin{equation}
V_{SRS}(T, F) = \int \mathrm{IF}^2(\mathbf{x}; T, F) \, dF(\mathbf{x}).
\end{equation}
Under RSS, the empirical distribution is modified by ranking-induced selection. For balanced RSS, each cycle contributes one observation from the $r$-th order statistic with density $f_{(r)}(\mathbf{x})$. Therefore, the asymptotic variance is
\begin{equation}
V_{RSS}(T, F) = \sum_{r=1}^{m} \pi_r \int \mathrm{IF}^2(\mathbf{x}; T, F) f_{(r)}(\mathbf{x}) \, d\mathbf{x}.
\end{equation}
To establish $V_{\mathrm{RSS}}(T, F) \leq V_{\mathrm{SRS}}(T, F)$, we use the variance decomposition identity. For balanced RSS with $\pi_r = 1/m$, write
\[
V_{\mathrm{SRS}}(T, F) = \int \mathrm{IF}^2(\mathbf{x}; T, F)\, dF(\mathbf{x}) = \frac{1}{m}\sum_{r=1}^{m}\int \mathrm{IF}^2(\mathbf{x}; T, F)\, f_{(r)}(\mathbf{x})\, d\mathbf{x} + R_m,
\]
where $R_m \geq 0$ captures the between-stratum variation in $\mathrm{IF}^2$ across the ranked strata. Since $\frac{1}{m}\sum_{r=1}^{m} f_{(r)}(\mathbf{x}) = f(\mathbf{x})$ for balanced RSS, the first term equals $V_{\mathrm{RSS}}(T, F)$, and the non-negative remainder $R_m$ yields
\begin{equation}
V_{\mathrm{RSS}}(T, F) = V_{\mathrm{SRS}}(T, F) - R_m \leq V_{\mathrm{SRS}}(T, F),
\end{equation}
provided the ranking induces non-degenerate stratification (i.e., the ranked strata are not identical to each other). This holds under perfect or near-perfect robust-compatible ranking.
\end{proof}
\begin{remark}\label{rem:nonsmooth_IF}
The influence function framework in Theorem~\ref{thm:asymp_robust_rss} applies directly to estimators with bounded, square-integrable influence functions, such as M-estimators with Huber or Tukey $\psi$-functions, and S-estimators with bounded $\rho$-functions. For the raw MCD and MVE estimators, whose influence functions are discontinuous and unbounded at the subset boundary~\citep{croux:haes:1999}, the theorem applies to their \emph{reweighted} versions, which possess bounded influence functions~\citep{maronna2019robust}. Alternatively, for the raw MCD and MVE, efficiency comparisons can be conducted via direct finite-sample variance ratios as in Section~5, bypassing the influence function apparatus.
\end{remark}

The results demonstrate that RSS can improve efficiency of robust estimators in multivariate settings if the ranking mechanism is robust. In particular, efficiency gain scales approximately with $\rho^2(\mathrm{RE}(1) - 1)$, where $\rho$ denotes ranking quality and $\mathrm{RE}(1)$ is the efficiency under perfect ranking. For instance, with $\rho = 0.7$ and $\mathrm{RE}(1) = 2$, required sample size for given precision can be reduced by ~33\% compared to SRS. In practice, set sizes of $m=4$ or $5$ are sufficient for meaningful variance reduction while preserving robustness.

\begin{theorem}[Breakdown Point Preservation in RSS]\label{thm:bdp_preservation}
%\label{thm:rss_breakdown}
Let $T$ be a robust estimator with finite-sample breakdown point $\epsilon^*$ under simple random sampling (SRS), and suppose RSS uses a robust-compatible ranking function with breakdown point $\epsilon^*_R \geq \epsilon^*$. If at least $\lceil m/2 \rceil$ of the $m$ ranking cycles contain fewer than $\lfloor m/2 \rfloor$ contaminated observations (i.e., a majority of cycles have a clean majority), then the RSS-based estimator $T_{RSS}$ satisfies
\[
\epsilon^*_{RSS} \geq \epsilon^* \cdot \left(1 - \frac{1}{2m} \right),
\]
Thus, $\varepsilon^*_{\mathrm{RSS}}$ approaches $\varepsilon^*$ from below as $m$ increases, with the gap bounded by $\varepsilon^*/(2m)$. For typical set sizes $m \geq 4$, the reduction is at most $12.5\%$ of $\varepsilon^*$, which is negligible in practice. Moreover, the bound is conservative: it assumes worst-case contamination concentration across cycles, which is unlikely under random allocation of contaminated observations to sets.
\end{theorem}

\begin{proof}
Consider worst-case contamination, where $\lfloor n\epsilon \rfloor$ of $n$ observations are replaced by outliers. Partitioning $n$ into $n/m$ cycles, if at least half the cycles have a clean majority, then each of these cycles contributes a clean observation even if some contamination exists. Since $\epsilon^*_R \geq \epsilon^*$, the ranking will select clean observations in most cycles. The effective contamination rate is thus bounded by
\[
\epsilon^*_{RSS} \geq \epsilon^* \cdot \left(1 - \frac{1}{2m} \right),
\]
concluding the proof.
\end{proof}
Thus, RSS preserves the breakdown point of the underlying estimator, provided the ranking is robust-compatible and contamination is not concentrated.
\begin{definition}[Efficiency–Robustness Frontier]
\label{def:frontier}
Let $\mathcal{T}$ denote a class of robust estimators, and $\mathcal{D}$ denote sampling designs (SRS, RSS with various rankings). The \emph{efficiency–robustness frontier} $\mathcal{F}$ is
\[
\mathcal{F} = \left\{ \left(V(T, D), \epsilon^*(T, D)\right) : T \in \mathcal{T}, D \in \mathcal{D} \right\},
\]
where $V(T, D)$ is asymptotic variance of $T$ under $D$, and $\epsilon^*(T, D)$ is its breakdown point.
\end{definition}
The preceding definition formalizes the notion of an \emph{efficiency–robustness frontier}, providing a unified framework to compare different robust estimators under various sampling designs. By mapping each estimator–design pair $(T, D)$ to its corresponding asymptotic variance and breakdown point, the frontier $\mathcal{F}$ succinctly captures the trade-offs between statistical efficiency and robustness. Theorem~\ref{thm:pareto_optimality} then establishes that, when combined with a robust-compatible RSS design, a high-breakdown estimator achieves a Pareto-optimal position on this frontier. In practical terms, this means that no alternative estimator–design combination can simultaneously improve efficiency and robustness without violating unbiasedness, highlighting the unique advantage of carefully integrating RSS with robust multivariate estimators. This perspective provides both theoretical guidance and practical insight for the design of sampling strategies that preserve estimator reliability while maximizing statistical precision.

\begin{theorem}[Pareto Optimality of RSS-Robust Systems]
\label{thm:pareto_optimality}
Suppose $T \in \mathcal{T}$ is a robust estimator with finite asymptotic variance and breakdown point, and $D \in \mathcal{D}$ is an RSS design with robust-compatible ranking. Then $(V(T,D), \epsilon^*(T,D))$ lies on the Pareto frontier $\mathcal{F}$. That is, no $(\tilde{T}, \tilde{D})$ exists such that
\[
V(\tilde{T}, \tilde{D}) < V(T,D) \quad \text{and} \quad \epsilon^*(\tilde{T}, \tilde{D}) > \epsilon^*(T,D),
\]
while preserving unbiasedness of $\tilde{T}$.
\end{theorem}

\begin{proof}
Suppose for contradiction that there exists $(T', D') \in \mathcal{T} \times \mathcal{D}$ with $V(T', D') < V(T, D)$ and $\varepsilon^*(T', D') > \varepsilon^*(T, D)$, where $T'$ is unbiased. We consider each inequality separately.

\emph{Variance bound.} By Theorem~\ref{thm:asymp_robust_rss}, under robust-compatible RSS, $V(T, D_{\mathrm{RSS}}) \leq V(T, D_{\mathrm{SRS}})$ for any fixed $T$. Hence, for the given estimator $T$, switching from SRS to RSS cannot increase variance. If $T'$ achieves $V(T', D') < V(T, D_{\mathrm{RSS}})$ while maintaining unbiasedness, then $T'$ under $D'$ has strictly lower variance than $T$ under RSS.

\emph{Breakdown constraint.} By Theorem~\ref{thm:bdp_preservation}, $\varepsilon^*(T, D_{\mathrm{RSS}}) \geq \varepsilon^*(1 - 1/(2m))$. The maximum achievable breakdown point for any affine equivariant estimator is $\lfloor(n - p + 1)/2\rfloor / n$ \citep{maronna2019robust}. If $T$ already achieves this upper bound (as MCD does), then $\varepsilon^*(T', D') > \varepsilon^*(T, D)$ requires $T'$ to exceed the theoretical maximum, a contradiction.

Thus, when $T$ achieves the maximal breakdown point and $D$ is robust-compatible RSS, no unbiased alternative can simultaneously improve both variance and breakdown point.
\end{proof}
This result justifies the use of RSS with robust estimators as an optimal trade-off between efficiency and robustness.


Theorem~\ref{thm:pareto_optimality} highlights that robust estimators paired with robust-compatible RSS designs achieve an optimal balance between efficiency and breakdown in the asymptotic regime. However, in practice, we often work with finite samples where asymptotic results provide only an approximation. To ensure the reliability and practical utility of these estimators, it is essential to examine their finite-sample behavior. The next subsection presents non-asymptotic results, offering concrete guarantees on variance, robustness, and efficiency for realistic sample sizes.
\subsection{Finite-Sample Properties of RSS-Robust Estimators}
\label{subsec:finite_sample}

While asymptotics describe long-run behavior, finite-sample properties are critical for practical applications. The following theorems provide non-asymptotic guarantees.

\begin{theorem}[Finite-Sample Efficiency Bound for the RSS Mean]
\label{thm:finite_eff}
Let $n = km$ denote the total number of measured RSS observations with set size $m$, $k$ cycles, and ranking correlation $\rho$. Then the relative efficiency of the RSS sample mean versus the SRS sample mean satisfies
\[
\text{RE}_n(RSS, SRS) \geq 1 + \frac{\rho^2}{n} \sum_{i=1}^{m} \frac{(E[X_{(i)}] - \mu)^2}{\sigma^2} + O(n^{-2}),
\]
with tightness for symmetric or exponential families.
\end{theorem}

\begin{proof}
Let $\bar{X}_{RSS}$ be the sample mean. A Taylor expansion around $E[X_{(i)}]$ gives
\[
\text{Var}(\bar{X}_{RSS}) = \frac{1}{n} \sum_{i=1}^{m} \text{Var}(X_{(i)}) + \text{Bias}^2,
\]
where $\text{Bias} = E[X_{(i)}] - \mu$, and $\text{Var}(X_{(i)}) \leq \sigma^2$. This yields
\[
\text{RE}_n = \frac{\text{Var}(\bar{X}_{SRS})}{\text{Var}(\bar{X}_{RSS})} \geq 1 + \frac{\rho^2}{n} \sum_{i=1}^{m} \frac{(E[X_{(i)}] - \mu)^2}{\sigma^2} + O(n^{-2}).
\]
\end{proof}
Theorem~\ref{thm:finite_eff} establishes how RSS improves estimator efficiency in finite samples, showing the role of ranking quality and set design in variance reduction. Building on this, Theorem~\ref{thm:finite_breakdown} examines the finite-sample robustness of RSS-based estimators, quantifying their breakdown point under contamination. Together, these results provide a practical understanding of the trade-offs between efficiency and robustness, highlighting the value of carefully designed RSS schemes for reliable multivariate estimation.

\begin{remark}
\label{rem:nonlinear_finite}
Theorem~\ref{thm:finite_eff} is established for the sample mean, which is a linear statistic. For nonlinear high-breakdown estimators such as MCD, S-estimators, and MM-estimators, closed-form finite-sample efficiency bounds are generally unavailable because these estimators lack explicit variance formulas in finite samples. For such estimators, we rely on the asymptotic results of Theorem~\ref{thm:asymp_robust_rss} and validate finite-sample performance empirically through the simulation studies in Section~\ref{sec:simulation}. Extending finite-sample theory to implicit, non-smooth estimators remains an important open problem.
\end{remark}

\begin{theorem}[Finite Sample Breakdown Point]
\label{thm:finite_breakdown}
Let $T_{RSS}$ be a robust estimator applied to an RSS sample of size $n$. The finite-sample breakdown point is
\[
\epsilon^*_{n,RSS} = \frac{1}{n} \max\left\{j : T_{RSS} \text{ remains bounded when } j \text{ observations are replaced arbitrarily} \right\},
\]
and if contamination is not concentrated in a majority of cycles,
\[
\epsilon^*_{n,RSS} \geq \epsilon^*_n,
\]
with equality in worst-case contamination.
\end{theorem}

\begin{proof}
Breakdown occurs when enough outliers force the estimator to diverge. If fewer than $\lfloor m/2 \rfloor$ contaminated units enter most cycles, bounded-influence rankings (Def.~\ref{def:robust_rss}) ensure clean selections dominate, giving
\[
\epsilon^*_{n,RSS} \geq \frac{1}{n} \left\lfloor \frac{n \cdot \epsilon^*_n}{1 + \delta} \right\rfloor,
\]
where $\delta = \max_{j} |\{i \in C_j : x_i \text{ is contaminated}\}|/m$ measures the maximum contamination fraction within any single RSS cycle $C_j$, and $\lfloor \cdot \rfloor$ denotes the floor function.
\end{proof}
In practice, the RSS breakdown point matches or exceeds that of SRS unless contamination is highly clustered.

The computational cost of RSS-enhanced robust estimators includes sorting ($O(km \log m)$), ranking ($O(km C_R)$), and estimation ($O(n C_T)$), where $C_R$ is the cost of ranking and $C_T$ the cost of the base robust estimator. In practice, RSS overhead (sorting and ranking) is typically negligible compared to estimation, especially for complex estimators like MCD or S-estimators, where $C_T$ grows quadratically or cubically with $n$.

In summary, RSS-based robust estimators achieve improved efficiency and maintain high breakdown points both asymptotically and in finite samples, provided the ranking function is robust-compatible. The next section presents empirical studies validating these theoretical results.

\section{Methodology and Implementation}
\label{sec:methodology}

This section develops the proposed RSS-enhanced robust estimation framework, integrating theoretical principles with practical implementation strategies. Rather than treating ranking functions, estimator-specific adaptations, and computational aspects in isolation, we present a unified methodology organized into four interrelated components: (i) the general framework, (ii) ranking strategies, (iii) estimator-specific implementations with tuning guidelines, and (iv) theoretical guarantees and performance considerations. This structure emphasizes both reproducibility and interpretability, ensuring that the methodology can be directly applied in practice.

\subsection{Unified Framework and Ranking Strategies}

At the core of our approach lies the combination of Ranked Set Sampling (RSS) with a robust estimator $T$. The procedure iteratively selects informative observations via a ranking function $R$, aggregates them into an RSS sample, and then applies $T$ for estimation. Algorithm~\ref{alg:general_rss_robust} summarizes this process.

\begin{algorithm}[H]
\caption{General RSS-Enhanced Robust Estimation}
\label{alg:general_rss_robust}
\begin{algorithmic}[1]
\State \textbf{Input:} Population $\mathcal{P}$, set size $m$, number of cycles $k$, robust estimator $T$, ranking function $R$
\State \textbf{Initialize:} RSS sample $\mathcal{S}_{RSS} = \emptyset$
\For{$i = 1$ to $k$}
    \State Draw $\{X_1, X_2, \ldots, X_m\}$ from $\mathcal{P}$
    \State Compute rankings $r_j = R(X_j)$ for $j=1,\ldots,m$
    \State Sort $X_{(1)}, \ldots, X_{(m)}$ by $r_j$
    \State Select $X_{(((i-1) \bmod m) + 1)}$ and add to $\mathcal{S}_{\mathrm{RSS}}$
\EndFor
\State Apply $T$ to $\mathcal{S}_{RSS}$ to obtain $\hat{\theta}$
\State \textbf{Return:} $\hat{\theta}$, diagnostics, confidence intervals
\end{algorithmic}
\end{algorithm}

The effectiveness of this framework depends critically on the choice of ranking function $R$. For multivariate settings, commonly used functions include:
\begin{itemize}
    \item \textbf{Robust Mahalanobis Distance:} $R(\mathbf{x}) = \mathbf{x}^T \hat{\boldsymbol{\Sigma}}^{-1} \mathbf{x}$
    \item \textbf{Projection Depth:} $R(\mathbf{x}) = \mathrm{med}_i |\mathbf{u}_i^T \mathbf{x}|$
    \item \textbf{Outlyingness:} $R(\mathbf{x}) = \max_{\|\mathbf{u}\|=1} \frac{|\mathbf{u}^T \mathbf{x} - \mathrm{med}|}{\mathrm{MAD}}$
    \item \textbf{$q$-Norm:} $R(\mathbf{x}) = \|\mathbf{x}\|_q$ (with $q=1,2,\infty$ being most common)
\end{itemize}

Each ranking criterion carries its own trade-off in terms of robustness, affine invariance, and computational burden. For example, Mahalanobis distance is efficient under normality but less robust in the presence of outliers, whereas projection depth and outlyingness provide greater robustness at higher computational cost. Selecting $R$ should therefore be guided by both data structure and contamination scenarios.

\subsection{Estimator-Specific Implementations and Tuning Guidelines}

Once a ranking function is chosen, different robust estimators can be integrated with RSS. Below we outline key cases.

%\paragraph{RSS-MCD.}
{\bf RSS-MCD.}

Here $R$ is typically the robust Mahalanobis distance. After constructing the RSS sample, the MCD estimator is applied with subset size
\[
h = \left\lfloor \frac{n + p + 1}{2} \right\rfloor,
\]
yielding breakdown point
\[
\epsilon^*_{\text{RSS-MCD}} \geq \min\left(\frac{n-h+1}{n}, \epsilon^*_R\right).
\]
Efficiency is improved by
\[
\text{ARE} = 1 + \rho^2 \cdot \kappa(m, p),
\]
where $\rho$ measures correlation between the ranking and the true centrality.

%\paragraph{RSS-S Estimator.}
{\bf RSS-S Estimator.}

Using projection depth as $R$, the S-estimator benefits from efficiency gain
\[
\text{ARE}(\text{RSS-S}, \text{S}) \geq 1 + \frac{c}{m} \rho^2,
\]
for constant $c$ depending on the $\rho$-function.

% \paragraph{RSS-Stahel–Donoho.}
{\bf RSS-Stahel–Donoho.}

Employing projection outlyingness as $R$, this approach reduces computational cost:
\[
\frac{C_{\text{RSS-SD}}}{C_{\text{SD}}} = O\left(\frac{1}{\sqrt{m}}\right),
\]
while maintaining robustness.

\vspace{0.5em}
In practice, successful implementation requires careful parameter tuning:
\begin{itemize}
    \item \textbf{Set Size ($m$):} Small values ($4 \leq m \leq 6$) often balance accuracy and cost.
    \item \textbf{Cycles ($k$):} Choose $k$ such that $n=km$ ensures desired precision.
    \item \textbf{Ranking Calibration:} Assess ranking correlation $\rho$ via simulations or clean subsets.
    \item \textbf{Affine Invariance:} Ensure $R$ is affine-invariant when using equivariant estimators like MCD.
    \item \textbf{Software:} Packages such as \texttt{robustbase} or \texttt{rrcov} can be adapted with ranking functions.
\end{itemize}

\subsection{Computational and Theoretical Considerations}

From a computational perspective, RSS introduces additional ranking overhead:
\[
\text{Cost}_{\text{RSS}} = O(km \log m) + O(km \cdot C_R),
\]
where $C_R$ is the cost of computing $R$. This is typically minor compared to the quadratic or higher complexity of high-breakdown estimators themselves.

Theoretically, RSS-based estimators inherit the robustness properties of their base estimators while offering improved efficiency. Under mild regularity conditions:
\begin{itemize}
    \item \textbf{Breakdown Point:} $\epsilon^*_{\text{RSS}} \geq \epsilon^*$ of the base estimator, provided $R$ is robust-compatible.
    \item \textbf{Efficiency Improvement:} Asymptotic variance decreases with stronger ranking correlation $\rho$ and moderate $m$.
    \item \textbf{Concentration Bounds:} Sub-Gaussian tails yield confidence intervals such as
    \[
    P(|\hat{\theta}_{\text{RSS}} - \theta| > t) \leq 2 \exp\left(-\frac{nt^2}{2(\sigma^2_{\text{RSS}} + \sigma^2_T)}\right).
    \]
\end{itemize}

Moreover, under contamination $F_\epsilon = (1-\epsilon)F_0 + \epsilon G$, the minimax risk satisfies
\[
\inf_{T_{\text{RSS}}} \sup_{F_\epsilon} E[\|T_{\text{RSS}} - \theta\|^2] = \Theta\left(\frac{1}{n} + \epsilon^2\right),
\]
establishing asymptotic optimality against adversarial noise.

\subsection{Performance Evaluation and Summary}

To assess empirical performance, we employ standard metrics:
\begin{itemize}
    \item \textbf{Breakdown Point (BP):} Smallest contamination fraction leading to breakdown.
    \item \textbf{Relative Efficiency (RE):} $\text{RE} = \text{MSE}(\hat{\theta}_{\text{SRS}})/\text{MSE}(\hat{\theta}_{\text{RSS}})$.
    \item \textbf{Convergence Rate (CR):} Typically $O(n^{-1/2})$, but improved with effective ranking.
\end{itemize}

In summary, the methodology integrates a general RSS-based framework, flexible ranking functions, estimator-specific implementations, and theoretical guarantees. Together these components ensure both robustness and efficiency, while remaining computationally feasible. The next section validates these claims through simulation studies.

\section{Simulation Study}
\label{sec:simulation}
We conducted a set of controlled simulations to evaluate how Ranked Set Sampling (RSS) improves the performance of robust estimators. The goal was to assess efficiency, robustness, inference quality, and computational cost under different contamination levels, ranking qualities, and sample sizes. All results are based entirely on simulated data, which allows precise control over the experimental conditions.

The simulation design was consistent across all experiments:
\begin{itemize}
    \item \textbf{Contamination levels}: $\epsilon \in \{0.05, 0.10, 0.20, 0.30, 0.40\}$, introduced as systematic outliers.
    \item \textbf{Ranking quality}: $\rho \in \{0.3, 0.5, 0.7, 0.9, 1.0\}$, where observed ranks are correlated with the true order via
    \[
    R_{\text{obs}} = \rho R_{\text{true}} + \sqrt{1-\rho^2}\,\varepsilon,
    \quad \varepsilon \sim N(0,1).
    \]
    \item \textbf{Sample sizes}: $n \in \{50, 100, 200\}$.
    \item \textbf{Set sizes}: $m \in \{3,5,7\}$ for balanced RSS designs.
    \item \textbf{Estimators}: sample mean, 20\% trimmed mean, Huber M-estimator (tuning constant $c_H = 1.345$), Tukey biweight (tuning constant $c_T = 4.685$), minimum covariance determinant (MCD), S-estimator, and MM-estimator.
    \item \textbf{Performance metrics}: relative efficiency (RE), relative bias, coverage probability, variance ratios, and computational cost (RSS/SRS runtime).
\end{itemize}

This framework provides a fair comparison across settings and highlights where RSS offers the most value.

\subsection{Univariate Simulation Results}
\label{subsec:univariate_simulation}

\textbf{Study 1: Ranking Quality and Efficiency.}
Table~\ref{tab:rank_eff} shows how efficiency depends on ranking quality. The empirical results closely matched the theoretical quadratic relationship, with only minor deviations. Efficiency gains were modest at poor ranking ($\rho=0.3$, RE $\approx 1.063$), but grew rapidly as ranking improved, reaching the theoretical maximum (RE $=2.618)$) under perfect ranking.

\begin{table}[htbp]
\centering
\caption{Ranking Quality vs. Relative Efficiency}\label{tab:rank_eff}
\begin{tabular}{cccccc}
\toprule
$\rho$ & SRS Var & RSS Var & Emp. RE & Theo. RE & Abs. Error \\
\midrule
0.30 & 0.1639 & 0.1542 & 1.063 & 1.188 & 0.125 \\
0.50 & 0.1650 & 0.1261 & 1.309 & 1.464 & 0.155 \\
0.70 & 0.1632 & 0.0906 & 1.801 & 1.993 & 0.192 \\
0.90 & 0.1646 & 0.0689 & 2.388 & 2.451 & 0.063 \\
1.00 & 0.1644 & 0.0628 & 2.6187 & 2.618 & 0.001 \\
\bottomrule
\end{tabular}
\end{table}

\textbf{Study 2: Robustness under Contamination.}
Table~\ref{tab:breakdown_analysis} reports estimator performance as contamination increases. The sample mean deteriorated quickly, with variance tripling by 40\% contamination. In contrast, robust estimators—especially Tukey and MM—remained stable. MCD and trimmed mean were more sensitive at higher contamination levels but still avoided full breakdown before 40\%.

\begin{table}[htbp]
\centering
\caption{Breakdown Analysis (Variance under Contamination)}
\label{tab:breakdown_analysis}
\begin{tabular}{lccccc}
\toprule
Estimator & 5\% & 10\% & 20\% & 30\% & 40\% \\
\midrule
Mean & 0.183 & 0.238 & 0.486 & 1.012 & 1.523 \\
Trimmed Mean & 0.161 & 0.180 & 0.275 & 0.441 & 0.889 \\
Huber M & 0.167 & 0.184 & 0.254 & 0.392 & 0.703 \\
Tukey Biweight & 0.164 & 0.176 & 0.241 & 0.355 & 0.621 \\
MCD & 0.171 & 0.187 & 0.515 & 1.018 & 1.448 \\
S-estimator & 0.170 & 0.186 & 0.266 & 0.422 & 0.782 \\
MM-estimator & 0.165 & 0.177 & 0.242 & 0.359 & 0.654 \\
\bottomrule
\end{tabular}
\end{table}

\textbf{Study 3: Efficiency Across Sample Sizes.}
Table~\ref{tab:rss_performance} shows how efficiency scaled with sample size. Trimmed mean and Huber M benefitted most, reaching nearly fourfold gains (RE $\approx 2.74$) at $n=200$. Even conservative estimators such as MCD and S showed steady, though smaller, improvements.

\begin{table}[htbp]
\centering
\caption{RSS vs. SRS: Relative Efficiency}
\label{tab:rss_performance}
\begin{tabular}{lcccc}
\toprule
Estimator & $n=50$ & $n=100$ & $n=200$ & Notes \\
\midrule
Mean & 1.55 & 2.21 & 2.68 & Gains scale with $n$ \\
Trimmed Mean & 1.64 & 2.72 & \textbf{2.74} & Strongest at $n=200$ \\
Huber M & 1.72 & 2.80 & \textbf{2.80} & Consistently strong \\
Tukey Biweight & 1.63 & 2.51 & 2.78 & Stable gains \\
MCD & 1.20 & 1.58 & 1.93 & Modest gains \\
S-estimator & 1.26 & 1.69 & 2.04 & Moderate gains \\
MM-estimator & 1.68 & 2.61 & 3.01 & Reliable across $n$ \\
\bottomrule
\end{tabular}
\end{table}

\textbf{Study 4: Small-Sample Inference.}
Table~\ref{tab:finite_sample} highlights coverage probabilities. RSS improved coverage to near-nominal levels for trimmed mean (from 0.83--0.85 under SRS to 0.95+ under RSS). MM-estimators tended to be conservative, with coverage near
$0.980$, while MCD lagged behind, sometimes below $0.75$.

\begin{table}[htbp]
\centering
\caption{Finite Sample Coverage and Variance Reduction}
\label{tab:finite_sample}
\begin{tabular}{lcccc}
\toprule
Estimator & $n$ & SRS Coverage & RSS Coverage & Var. Ratio \\
\midrule
Trimmed Mean & 50 & 0.857 & \textbf{0.956} & 2.47 \\
Trimmed Mean & 100 & 0.832 & \textbf{0.971} & 2.77 \\
Trimmed Mean & 200 & 0.830 & \textbf{0.970} & 2.68 \\
MM-estimator & 50 & 1.000 & 0.980 & 2.55 \\
MM-estimator & 100 & 1.000 & 0.961 & 2.77 \\
MCD & 50 & 0.743 & 0.794 & 1.40 \\
MCD & 100 & 0.659 & 0.732 & 1.58 \\
\bottomrule
\end{tabular}
\end{table}

\textbf{Study 5: Computational Cost.}
Finally, Table~\ref{tab:computational} compares runtime. For simple estimators (mean, trimmed), RSS was slightly faster than SRS because of reduced variability. For complex estimators like MCD, runtime was essentially the same. Overall, efficiency gains came with no computational penalty.

\begin{table}[htbp]
\centering
\caption{Computational Overhead (RSS/SRS Time Ratio)}
\label{tab:computational}
\begin{tabular}{lcccc}
\toprule
Estimator & $n=50$ & $n=100$ & $n=200$ & $n=500$ \\
\midrule
Mean & 0.72 & 0.68 & 0.62 & 0.59 \\
Trimmed Mean & 0.81 & 0.79 & 0.80 & 0.81 \\
Huber M & 0.92 & 0.91 & 0.89 & 0.90 \\
Tukey Biweight & 0.97 & 0.97 & 0.98 & 0.98 \\
MCD & 0.99 & 1.01 & 0.97 & 0.95 \\
\bottomrule
\end{tabular}
\end{table}

\bigskip
\noindent
\textbf{Note on computational timing.} The RSS/SRS time ratios below 1.0 for simple estimators (mean, trimmed mean) in Table~\ref{tab:computational} reflect the fact that RSS, by reducing estimator variance, requires fewer Monte Carlo replications to achieve the same precision target in variance estimation. When computational cost is measured per single estimate (rather than per precision-matched experiment), the RSS procedure is slightly more expensive due to the ranking step. For complex iterative estimators (MCD, S-estimator), the ratios are near 1.0, confirming that the ranking overhead is negligible relative to the estimation cost. We report both per-estimate and precision-matched timing to avoid ambiguity.

In summary, these simulations show that RSS meaningfully boosts efficiency (up to $\approx 2.8$ max),
helps robust estimators maintain stability under contamination, improves small-sample inference, and adds little to no computational burden. Among all methods, trimmed mean and Huber M stood out as the most reliable beneficiaries of RSS.

% \section{Multivariate Simulation Study}
% \label{sec:multivariate_simulation}
\subsection{Multivariate Simulation Study}
\label{subsec:multivariate_simulation}

This section presents multivariate RSS simulations across dimensions $p \in \{4, 10, 50\}$, using 1000 Monte Carlo replications per configuration. The ranking function is the robust Mahalanobis distance based on the trimmed covariance matrix (Section~\ref{sec:methodology}). Performance was evaluated across contamination levels, ranking qualities, and sample and set sizes, focusing on relative efficiency (RE), design effects, and the impact of dimensionality on RSS gains.

\subsubsection{Ranking Quality Across Dimensions}

Tables~\ref{tab:mv_rank_eff_p4}--\ref{tab:mv_rank_eff_p50} report the relationship between ranking quality $\rho$ and empirical relative efficiency for $p = 4$, $p = 10$, and $p = 50$, respectively. Unlike the univariate setting (Table~\ref{tab:rank_eff}), where RSS achieved RE values up to 2.6, the multivariate empirical RE values cluster near 1.0 across all dimensions and ranking qualities. This reflects the difficulty of achieving effective multivariate ranking: depth-based ranking in $\mathbb{R}^p$ captures less distributional information per observation than univariate ranking, reducing the variance-reduction benefit of RSS. The theoretical bounds, derived under univariate assumptions, are correspondingly conservative in the multivariate setting.

\begin{table}[htbp]
\centering
\caption{Ranking Quality vs.\ Statistical Efficiency ($p = 4$, 1000 replications)}
\label{tab:mv_rank_eff_p4}
\begin{tabular}{cccc}
\toprule
$\rho$ & Emp.\ RE & Theo.\ RE & Abs.\ Error \\
\midrule
0.30 & 0.929 & 1.180 & 0.251 \\
0.50 & 0.985 & 1.500 & 0.515 \\
0.70 & 0.945 & 1.980 & 1.035 \\
0.90 & 0.948 & 2.620 & 1.672 \\
1.00 & 0.969 & 3.000 & 2.031 \\
\midrule
\multicolumn{3}{l}{Mean Absolute Error} & 1.101 \\
\bottomrule
\end{tabular}
\end{table}

\begin{table}[htbp]
\centering
\caption{Ranking Quality vs.\ Statistical Efficiency ($p = 10$, 1000 replications)}
\label{tab:mv_rank_eff_p10}
\begin{tabular}{cccc}
\toprule
$\rho$ & Emp.\ RE & Theo.\ RE & Abs.\ Error \\
\midrule
0.30 & 0.911 & 1.180 & 0.269 \\
0.50 & 0.949 & 1.500 & 0.551 \\
0.70 & 0.926 & 1.980 & 1.054 \\
0.90 & 0.930 & 2.620 & 1.690 \\
1.00 & 0.968 & 3.000 & 2.032 \\
\midrule
\multicolumn{3}{l}{Mean Absolute Error} & 1.119 \\
\bottomrule
\end{tabular}
\end{table}

\begin{table}[htbp]
\centering
\caption{Ranking Quality vs.\ Statistical Efficiency ($p = 50$, 1000 replications)}
\label{tab:mv_rank_eff_p50}
\begin{tabular}{cccc}
\toprule
$\rho$ & Emp.\ RE & Theo.\ RE & Abs.\ Error \\
\midrule
0.30 & 0.946 & 1.180 & 0.234 \\
0.50 & 0.924 & 1.500 & 0.576 \\
0.70 & 0.942 & 1.980 & 1.038 \\
0.90 & 0.936 & 2.620 & 1.684 \\
1.00 & 0.931 & 3.000 & 2.069 \\
\midrule
\multicolumn{3}{l}{Mean Absolute Error} & 1.120 \\
\bottomrule
\end{tabular}
\end{table}
%==============paragraph

\subsubsection{Estimator Performance Under Contamination}

Tables~\ref{tab:contam_p4_new}--\ref{tab:contam_p50_new} report relative efficiency for four estimators under increasing contamination at $p = 4$, $p = 10$, and $p = 50$. The results reveal a nuanced pattern: RSS gains are estimator- and contamination-dependent rather than universal.

At $p = 4$ (Table~\ref{tab:contam_p4_new}), the trimmed mean benefits most from RSS, reaching RE $= 1.317$ at 20\% contamination. Huber~M shows modest gains at 30\% contamination (RE $= 1.089$), while the sample mean and Fast~MCD remain near RE $= 1.0$. At $p = 10$ (Table~\ref{tab:contam_p10_new}), gains are smaller but still present: the trimmed mean achieves RE $= 1.104$ at 20\% contamination, and Fast~MCD shows a notable gain (RE $= 1.347$) at 10\% contamination. At $p = 50$ (Table~\ref{tab:contam_p50_new}), the trimmed mean continues to benefit from RSS (RE $= 1.284$ at 20\%), and Fast~MCD shows consistent gains across contamination levels (RE $= 1.05$--$1.15$). These results indicate that RSS gains persist in higher dimensions for estimators whose ranking function aligns well with the contamination structure, though the magnitude of improvement diminishes.

\begin{table}[htbp]
\centering
\caption{Estimator Performance Under Contamination ($p = 4$, 1000 replications)}
\label{tab:contam_p4_new}
\begin{tabular}{llcc}
\toprule
Contam. & Estimator & Rel.\ Eff. & SRS MSE \\
\midrule
5\% & Sample Mean & 1.045 & 1.386 \\
& Trimmed Mean & 0.962 & 0.111 \\
& Huber M & 0.949 & 0.095 \\
& Fast MCD & 0.985 & 0.219 \\
\midrule
10\% & Sample Mean & 0.992 & 4.584 \\
& Trimmed Mean & 1.112 & 0.230 \\
& Huber M & 0.983 & 0.150 \\
& Fast MCD & 0.922 & 0.203 \\
\midrule
20\% & Sample Mean & 0.984 & 16.648 \\
& Trimmed Mean & \textbf{1.317} & 2.930 \\
& Huber M & 1.000 & 0.438 \\
& Fast MCD & 1.025 & 0.254 \\
\midrule
30\% & Sample Mean & 1.010 & 37.428 \\
& Trimmed Mean & 1.001 & 17.884 \\
& Huber M & 1.089 & 1.431 \\
& Fast MCD & 1.056 & 3.735 \\
\bottomrule
\end{tabular}
\end{table}

\begin{table}[htbp]
\centering
\caption{Estimator Performance Under Contamination ($p = 10$, 1000 replications)}
\label{tab:contam_p10_new}
\begin{tabular}{llcc}
\toprule
Contam. & Estimator & Rel.\ Eff. & SRS MSE \\
\midrule
5\% & Sample Mean & 1.020 & 3.372 \\
& Trimmed Mean & 0.955 & 0.278 \\
& Huber M & 0.946 & 0.220 \\
& Fast MCD & 0.915 & 0.400 \\
\midrule
10\% & Sample Mean & 1.026 & 11.301 \\
& Trimmed Mean & 1.078 & 0.581 \\
& Huber M & 0.969 & 0.345 \\
& Fast MCD & \textbf{1.347} & 0.744 \\
\midrule
20\% & Sample Mean & 0.992 & 42.769 \\
& Trimmed Mean & 1.104 & 6.407 \\
& Huber M & 1.021 & 1.001 \\
& Fast MCD & 0.927 & 9.764 \\
\midrule
30\% & Sample Mean & 0.985 & 92.435 \\
& Trimmed Mean & 1.047 & 43.560 \\
& Huber M & 1.037 & 3.152 \\
& Fast MCD & 0.975 & 60.543 \\
\bottomrule
\end{tabular}
\end{table}

\begin{table}[htbp]
\centering
\caption{Estimator Performance Under Contamination ($p = 50$, 1000 replications)}
\label{tab:contam_p50_new}
\begin{tabular}{llcc}
\toprule
Contam. & Estimator & Rel.\ Eff. & SRS MSE \\
\midrule
5\% & Sample Mean & 1.077 & 17.770 \\
& Trimmed Mean & 0.996 & 1.416 \\
& Huber M & 0.960 & 1.064 \\
& Fast MCD & 1.154 & 17.387 \\
\midrule
10\% & Sample Mean & 1.034 & 58.551 \\
& Trimmed Mean & 1.082 & 2.953 \\
& Huber M & 0.987 & 1.669 \\
& Fast MCD & 1.072 & 59.021 \\
\midrule
20\% & Sample Mean & 1.026 & 212.047 \\
& Trimmed Mean & \textbf{1.284} & 36.388 \\
& Huber M & 0.988 & 4.728 \\
& Fast MCD & 1.052 & 220.516 \\
\midrule
30\% & Sample Mean & 0.993 & 467.804 \\
& Trimmed Mean & 1.012 & 215.605 \\
& Huber M & 1.006 & 15.392 \\
& Fast MCD & 1.031 & 480.266 \\
\bottomrule
\end{tabular}
\end{table}
%==================================== paragraph

\subsubsection{Design Effects Across Dimensions}

Table~\ref{tab:design_effects_all} reports the relative efficiency of the trimmed mean under 10\% contamination for different $(n, m)$ configurations across the three dimensions. Two consistent patterns emerge: (i)~larger set sizes ($m = 7$) generally yield higher RE than smaller ones ($m = 3$), confirming that more ranking information improves RSS performance; and (ii)~smaller total sample sizes ($n = 30$) benefit more from RSS than larger ones ($n = 60$), consistent with the principle that RSS is most valuable when observations are expensive. At $p = 4$, the best configuration is $(n = 30, m = 7)$ with RE $= 1.384$. At $p = 10$, the same configuration achieves RE $= 1.422$. At $p = 50$, gains are more modest but still favor $m = 7$ (RE $= 1.246$).

\begin{table}[htbp]
\centering
\caption{Design Effects: Relative Efficiency of Trimmed Mean ($\epsilon = 0.10$, $\rho = 0.9$, 1000 replications)}
\label{tab:design_effects_all}
\begin{tabular}{cc|ccc}
\toprule
$n$ & $m$ & $p = 4$ & $p = 10$ & $p = 50$ \\
\midrule
30 & 3 & 1.094 & 1.226 & 1.082 \\
30 & 5 & 1.300 & 1.256 & 1.097 \\
30 & 7 & \textbf{1.384} & \textbf{1.422} & \textbf{1.246} \\
60 & 3 & 1.029 & 0.922 & 0.956 \\
60 & 5 & 1.102 & 1.031 & 1.028 \\
60 & 7 & 1.287 & 1.105 & 1.167 \\
\bottomrule
\end{tabular}
\end{table}
%=============paragraph

\subsubsection{Cross-Dimensional Summary}

Table~\ref{tab:cross_dim_summary} provides a unified comparison of RSS performance across dimensions at 10\% contamination with $\rho = 0.9$. The results show that RSS gains are modest in multivariate settings (RE between 0.80 and 1.09) compared to the univariate case, but remain consistently above or near 1.0 for most estimator--dimension combinations. The trimmed mean and Huber~M show the most reliable gains, while Fast~MCD exhibits more variable performance.

\begin{table}[htbp]
\centering
\caption{Cross-Dimensional Relative Efficiency Summary ($\epsilon = 0.10$, $\rho = 0.9$, $n = 60$, $m = 5$, 1000 replications)}
\label{tab:cross_dim_summary}
\begin{tabular}{lcccc}
\toprule
Dimension & Sample Mean & Trimmed Mean & Huber M & Fast MCD \\
\midrule
$p = 4$ & 1.062 & 0.968 & 1.031 & 0.936 \\
$p = 10$ & 1.017 & 0.978 & 1.001 & 0.799 \\
$p = 50$ & 1.093 & 1.048 & 1.017 & 1.074 \\
\bottomrule
\end{tabular}
\end{table}

Several observations emerge from the multivariate experiments. First, RSS gains in multivariate settings are substantially more modest than in univariate settings, primarily because depth-based ranking in $\mathbb{R}^p$ captures less distributional information per observation than univariate ranking. Second, the benefits are estimator- and contamination-dependent: the trimmed mean benefits most under moderate-to-high contamination (20--30\%), while the sample mean and Fast~MCD show smaller, more variable gains. Third, larger set sizes ($m = 7$) consistently outperform smaller ones across all dimensions, confirming that richer ranking information improves RSS performance even in high dimensions. Finally, the theoretical efficiency bounds derived under univariate assumptions are conservative in the multivariate setting, suggesting that tighter multivariate-specific bounds are a valuable direction for future work.



% \section{Real Data Application}
% \label{sec:real_data}
\subsection{Illustrative Applications on Benchmark Datasets}\label{sec:realdata}
To complement the simulation study and highlight the practical value of RSS, we applied the method to two well-known benchmark datasets: the IRIS botanical dataset ($n=150, p=4$) and a subset of the Wine Quality dataset ($n=300, p=4$).
These datasets were chosen because they provide clear and interpretable auxiliary ranking variables—Mahalanobis distance for IRIS and recorded quality scores for Wine—which make it easy to see how RSS leverages side information in practice.
For each dataset, we considered two estimators, the sample mean and the 20\% trimmed mean, and evaluated performance under 10\% contamination across varying sample sizes.
Relative efficiency (RE), defined as $\text{Var}_{\text{SRS}}/\text{Var}_{\text{RSS}}$, was used as the primary performance measure, averaged over 100 Monte Carlo replications.

%\paragraph{IRIS Dataset}
{\bf IRIS Dataset.}
We first evaluated RSS performance on the classic IRIS dataset ($n=150, p=4$), using Mahalanobis distance as the ranking variable.
Subsamples of size $n=\{20,40,60\}$ with set size $m=4$ were analyzed under 10\% contamination.
Table~\ref{tab:iris_results} shows that the trimmed mean benefited consistently from RSS, with efficiency gains of about 10\% across all sample sizes.
The sample mean improved slightly at smaller $n$ but fell below one at $n=60$ (RE=0.684), indicating that contamination effects combined with imperfect ranking can offset RSS advantages in larger samples.
Overall, the results suggest that RSS paired with robust estimators such as the trimmed mean yields the most reliable improvements.

\begin{table}[htbp]
\centering
\caption{IRIS Dataset: Relative Efficiency (100 replications, 10\% contamination)}
\label{tab:iris_results}
\begin{tabular}{lccc}
\hline
Estimator & $n=20$ & $n=40$ & $n=60$ \\
\hline
Sample Mean   & 1.035 & 0.954 & 0.684 \\
Trimmed Mean  & 1.121 & 1.028 & 1.126 \\
\hline
\end{tabular}
\end{table}

%\paragraph{Wine Quality Dataset}

{\bf Wine Quality Dataset.}
We next considered the Wine Quality dataset ($n=300, p=4$), using the recorded quality scores as auxiliary ranking information.
Here, sample sizes $n=\{20,40,60\}$ with set size $m=4$ were tested under 10\% contamination.
As shown in Table~\ref{tab:wine_results}, RSS produced clear efficiency gains in most scenarios.
At smaller samples ($n=20,40$), both the sample mean and trimmed mean improved, with trimmed mean reaching a peak efficiency of RE=1.320 at $n=40$.
At $n=60$, the sample mean continued to benefit (RE=1.254), while the trimmed mean dropped slightly below one (RE=1.120).
This pattern suggests that the usefulness of auxiliary rankings may diminish as sample size increases, but robust estimators can still leverage RSS effectively when the ranking variable is strongly informative.

\begin{table}[htbp]
\centering
\caption{Wine Quality Dataset: Relative Efficiency (100 replications, 10\% contamination)}
\label{tab:wine_results}
\begin{tabular}{lccc}
\hline
Estimator & $n=20$ & $n=40$ & $n=60$ \\
\hline
Sample Mean   & 1.264 & 1.116 & 1.254 \\
Trimmed Mean  & 1.167 & 1.320 & 1.120 \\
\hline
\end{tabular}
\end{table}

%\paragraph{Cross-Dataset Takeaways}
{\bf Cross-Dataset Takeaways.}
Taken together, these two real-data studies show that RSS can provide meaningful efficiency gains, particularly when auxiliary ranking variables are well aligned with the data structure.
On average, improvements were modest but consistent—about 8\% for IRIS and 17\% for Wine—demonstrating that even simple ranking schemes can translate into tangible benefits.
At the same time, the results also underscore that RSS is not universally superior to simple random sampling: efficiency occasionally fell below one, especially at larger sample sizes or when the ranking signal weakened.
This balanced evidence suggests that RSS is best viewed as a complementary tool that enhances robust estimators in practical settings, rather than a guaranteed substitute for random sampling.


% \paragraph{Air Quality Dataset}
{\bf Air Quality Dataset.}
For $n=\{40,80,120\}$ with $m=6$, Table~\ref{tab:air_quality_results} demonstrates consistent RSS efficiency improvements across all estimators. The analysis achieved strong correlations between the Air Quality Index (AQI) ranking variable and study variables: AQI vs Ozone ($r=0.964$), AQI vs Wind ($r=-0.934$), AQI vs Temperature ($r=0.868$), and AQI vs Solar Radiation ($r=0.770$). Sample Mean estimator recorded the highest efficiency RE=1.817 at $n=40$, while maintaining substantial improvements of 55.8-60.5\% for larger sample sizes. Trimmed Mean demonstrated remarkable stability (RE=1.571-1.738), and Huber M-estimator showed consistent performance (RE=1.578-1.743). Fast MCD, while exhibiting lower relative efficiency (RE=1.135-1.432), still provided meaningful improvements over SRS. The average relative efficiency of 1.552 across all estimators represents a 55.2\% improvement in estimation precision, with efficiency gains ranging from 13.5\% to 81.7\%.

\begin{table}[htbp]
\centering
\caption{Air Quality Dataset: Relative Efficiency Results}
\label{tab:air_quality_results}
\begin{tabular}{lccc}
\hline
Estimator & $n=40$ & $n=80$ & $n=120$ \\
\hline
Sample Mean & \textbf{1.817} & 1.605 & 1.558 \\
Trimmed Mean & 1.726 & \textbf{1.738} & 1.571 \\
Huber M & 1.743 & 1.578 & 1.585 \\
Fast MCD & 1.432 & 1.141 & 1.135 \\
\hline
\textit{Average} & \textit{1.680} & \textit{1.516} & \textit{1.462} \\
\hline
\end{tabular}
\end{table}

\begin{table}[htbp]
\centering
\caption{Air Quality Dataset: Variance Comparison (SRS vs RSS)}
\label{tab:air_quality_variance}
\begin{tabular}{lccc}
\hline
Estimator & $n=40$ & $n=80$ & $n=120$ \\
\hline
Sample Mean & 42.402/23.336 & 18.628/11.607 & 12.260/7.871 \\
Trimmed Mean & 43.907/25.446 & 21.162/12.175 & 12.153/7.735 \\
Huber M & 45.496/26.106 & 19.684/12.472 & 13.365/8.431 \\
Fast MCD & 96.740/67.563 & 42.385/37.134 & 28.343/24.967 \\
\hline
\multicolumn{4}{l}{\footnotesize Note: Values shown as SRS/RSS variance}
\end{tabular}
\end{table}

The variance reduction analysis in Table~\ref{tab:air_quality_variance} reveals substantial improvements in estimation precision. Sample Mean achieved variance reductions of 45\%, 38\%, and 36\% for $n=\{40,80,120\}$ respectively. The strong correlation structure ($|r| > 0.77$ for all variables) enables effective RSS ranking, resulting in consistent efficiency gains across all sample sizes. Unlike previous studies reporting sporadic high efficiency values, these results demonstrate the reliable and theoretically sound performance of RSS with robust estimators in environmental monitoring applications.

% \paragraph{Key Findings}
{\bf Key Findings.}
Our analysis highlights several important insights. First, RSS consistently provides modest but reliable efficiency gains when applied correctly, averaging about 8.2\% across the datasets we examined. Second, robust estimators—especially the trimmed mean—tend to benefit the most from RSS, showing steady improvements in the 10–13\% range. Third, the quality of auxiliary information plays a crucial role: for example, the Wine dataset enjoyed a strong 17\% average gain, whereas the IRIS dataset saw more modest improvements. Finally, sample size influences RSS performance, with smaller samples generally showing stronger relative efficiency gains. Overall, these findings suggest that RSS offers meaningful, practical improvements over simple random sampling, particularly in the presence of data contamination.

%\paragraph{Cross-Dataset Comparison}
{\bf Cross-Dataset Comparison.}
Table~\ref{tab:cross_domain} provides a snapshot of the relative efficiency results across the two datasets. The Wine dataset stands out for its consistent performance, while IRIS shows more variable outcomes. The higher average RE for Wine (1.173) compared to IRIS (0.991) reinforces the importance of carefully selecting auxiliary ranking variables when implementing RSS.

\begin{table}[htbp]
\centering
\caption{Cross-Domain Relative Efficiency Summary}
\label{tab:cross_domain}
\begin{tabular}{l c c c}
\hline
Estimator & IRIS Average & Wine Average & Overall Average \\
\hline
Sample Mean & 0.891 & 1.211 & 1.051 \\
Trimmed Mean & 1.092 & 1.136 & 1.114 \\
\hline
Dataset Average & 0.991 & 1.173 & 1.082 \\
\hline
\end{tabular}
\end{table}

In short, these real-data applications show that RSS combined with robust estimators can deliver meaningful efficiency gains. The degree of improvement depends heavily on both the choice of estimator and the quality of the auxiliary ranking information, highlighting RSS as a practical and effective complement to standard random sampling approaches.

\section{Conclusion}\label{sec:conclusion}

This paper has presented a unified framework for integrating Ranked Set Sampling (RSS) with high-breakdown robust estimators, including MCD, MVE, S-estimators, and the Stahel-Donoho estimator. The main theoretical contributions are: (i)~a closed-form characterization of RSS efficiency under imperfect ranking (Theorem~\ref{thm:rss_efficiency}), with an exact expression and a distribution-free lower bound; (ii)~asymptotic variance reduction guarantees for robust estimators applied to RSS samples (Theorem~\ref{thm:asymp_robust_rss}), valid for estimators with bounded influence functions; and (iii)~breakdown point preservation results showing that robust-compatible ranking maintains the robustness of the underlying estimator (Theorem~\ref{thm:bdp_preservation}).

Simulation studies in univariate settings confirmed that RSS yields substantial efficiency gains (relative efficiency up to 2.8) while preserving robustness under contamination levels up to 30\%. Multivariate simulations across $p \in \{4, 10, 50\}$ with 1000 replications revealed a more nuanced picture: RSS provides meaningful gains for specific estimator--contamination combinations (particularly the trimmed mean at moderate contamination, with RE up to 1.32 at $p = 4$ and 1.28 at $p = 50$), while gains are modest overall and diminish with increasing dimensionality. Illustrative applications on the IRIS, Wine Quality, and Air Quality benchmark datasets demonstrated consistent, though modest, improvements in estimation precision.

Several limitations should be noted. First, the finite-sample efficiency bounds (Theorem~\ref{thm:finite_eff}) are established only for linear statistics; extending these to nonlinear high-breakdown estimators remains an open problem. Second, while the multivariate simulations now cover $p$ up to 50, the regime $p \gg n$ (where $p$ exceeds the sample size) remains unexplored and poses additional challenges for depth-based ranking. Third, the theoretical results assume that contamination is independent of the ranking mechanism; when outliers systematically affect ranking, additional safeguards may be needed. Fourth, the theoretical efficiency bounds derived under univariate assumptions are conservative in multivariate settings, as demonstrated by the gap between theoretical and empirical RE values in Tables~\ref{tab:mv_rank_eff_p4}--\ref{tab:mv_rank_eff_p50}; developing tighter multivariate-specific bounds remains open.

Future work should address: (a)~extension to the high-dimensional regime $p > n$, with comparison against recent scalable robust methods such as DetMCD \citep{Hubert2012} and cellwise MCD \citep{Raymaekers2023}; (b)~development of multivariate-specific efficiency bounds that account for the reduced informativeness of depth-based ranking in high dimensions; (c)~adaptive selection of the set size $m$ and ranking function based on data characteristics and dimensionality; and (d)~theoretical analysis of the interaction between ranking errors and breakdown point in adversarial contamination models.


\section*{Funding}
No funding was received for conducting this study.

\section*{Competing Interests}
The authors declare that they have no competing interests.


%\section{BibTeX}\label{bibtex}

%We encourage you to use BibTeX. If you have, please feel free to use the
%package natbib with any bibliography style you're comfortable with. The
%.bst file agsm has been included here for your convenience.


\bibliography{bibliography.bib}


\end{document}
